\documentclass[11pt,a4paper,oneside]{book}

%%%% PACKAGES WITH OPTIONS
\usepackage[english]{babel} % Language-specific hyphenation/splitting
\usepackage[small,bf,centerlast]{caption}[=v1] % Customize appearance of fig/table captions
\usepackage[Sonny]{fncychap} % Styles for chapter headers: Bjarne, Bjornstrup, Conny, Glenn, Lenny, Rejne and Sonny
\usepackage[T1]{fontenc} % Improved font encoding for better output, especially for accented characters
% \usepackage[inner=2.5cm,outer=2.5cm,vmargin=2.5cm]{geometry} % Set page margins
\usepackage[final]{graphicx} % Handles inclusion of images
\usepackage[UTF8]{inputenc} % Support for Latin-1 encoding (UTF-8 is generally preferred)\\
\usepackage[version=4]{mhchem} % Chemistry and isotopes with \ce{} (other option: 'chemformula' \ch{})
\usepackage[square,sort&compress,numbers]{natbib} % For formatted bibliographies with numeric citation styles
\usepackage[final]{pdfpages} % Inclusion of external PDF files
\usepackage{setspace} % Controls line spacing; "onehalfspacing" adds readability
\usepackage[separate-uncertainty=true]{siunitx} % Fancy display of SI units, can show uncertainty (e.g. \SI{1.54(6)}{\tesla} --> 1.54 ± 0.06 T)
\usepackage[dvipsnames]{xcolor} % Advanced color handling with named colors

%%%% PACKAGES WITHOUT OPTIONS
\usepackage{a4wide} % Adjusts margins for A4 paper (could conflict with "geometry")
%\usepackage{alphalph,cite} % Alphabetic-style labels
\usepackage{amsfonts}
\usepackage{amsmath} % Advanced math formatting and features
\usepackage{amssymb} % Additional math symbols
\usepackage{amsthm} % Environments like "proof"
\usepackage{bm} % Bold math symbols
\usepackage{cancel} % Strikethrough diagonal in math mode
\usepackage{changepage} % Enables adjustable margins within the document
\usepackage{color} % Basic color support (often superseded by "xcolor")
\usepackage{enumitem}
\usepackage{float} % Provides option [H] to floats
\usepackage[backref=none]{hyperref} % Creates hyperlinks and metadata in PDFs. Set backref=page to add back-references in bibliography.
\usepackage{listings} % Formats code snippets for programming examples
\usepackage{lmodern} % Modern font family with better scalability and appearance
\usepackage{longtable} % Tables that can span multiple pages
\usepackage{lscape} % Rotate content to landscape orientation
\usepackage{makecell} % Line breaks in table cells
\usepackage{makeidx} % Creates an index for the document
\usepackage[framemethod=TikZ]{mdframed} % Creates framed content with custom styles
\usepackage{microtype} % Enhances typographic appearance (e.g., kerning and justification)
\usepackage{nonfloat} % Prevents floating figures/tables for specific environments
\usepackage{pagecolor} % Sets background color of pages
\usepackage{physics2} % Simplifies common physics expressions and notation
\usepackage{ragged2e} % Provides \justifying command
\usepackage{soul} % For text decoration, e.g., underline (\ul{}), strikethrough (\st{}), highlighting...
% \usepackage{subcaption} % Subfigure captions
\usepackage{tikz}
\usepackage{titlesec} % Customizes the format and spacing of section titles (e.g. \setcounter{secnumdepth} to make subsubsections numbered in chapters)
\usepackage{tocbibind} % Adds bibliography and table of contents to ToC
\usepackage{tocloft} % Customizes the layout of the table of contents
\usepackage{upgreek} % Upright Greek letters
\usepackage{wasysym} % Many symbols
\usepackage{xfrac} % For diagonal fractions in text mode with \sfrac{}{}
\usepackage{xspace} % Automatically manages spaces after commands
\usepackage{xurl} % Better URL handling with line breaks

%%%% PACKAGES THAT MUST BE LOADED AFTER OTHERS
\usepackage[all]{hypcap} % [Requires "hyperref"] Ensures links to figures focus on the image, not the caption
\usepackage[capitalise]{cleveref} % [Requires "hyperref"] Intelligent cross-referencing

%%%% CREATE CUSTOM COMMANDS
% \definecolor{figcolor}{HTML}{F2F2FF}
\definecolor{figcolor}{HTML}{F2F2F2}
\definecolor{ugentblue}{RGB}{10,30,96}
\definecolor{lightblue}{HTML}{1E90FF} % C0
\definecolor{lightred}{HTML}{D62728} % C1
\definecolor{mplgrey}{gray}{0.5} % "grey" in matplotlib

\newcommand{\colorcode}{(Color code on page \pageref{colorcode}.)\xspace}

%% Index labels.
% To avoid multiply-defined labels, but allow multiple pages in the index, use e.g. \idx[nolabel]{...} after the first \idx{...}. 
\newcommand{\indexlabel}[2][]{\index{#2}\ifthenelse{\equal{#1}{nolabel}}{}{\label{#2}\label{#2s}}}
\newcommand{\idx}[2][]{\emph{#2}\indexlabel[#1]{#2}} %no xspace
\newcommand{\xlabel}[2][]{#2\indexlabel[#1]{#2}}
%% Index references.
\newcommand{\xref}[1]{\hyperref[#1]{#1}\xspace}
\newcommand{\xrefs}[2]{\hyperref[#2]{#1#2}\xspace} %% meervoud etc, add {s}
\newcommand{\link}[2]{\hyperref[#1]{#2}\xspace}

\newcommand{\jonathanskip}{\smallskip}
\newcommand{\numberthis}{\addtocounter{equation}{1}\tag{\theequation}} % Numbers a single line in a no-numbering multiline equation* or align*
\newcommand{\numeq}[2]{\stackrel{\scriptscriptstyle(\mkern-1.5mu#1\mkern-1.5mu)}{#2}} % Number above mathematical symbol (usage: \numeq{<eq_number>}{<symbol>})

%\newcommand{\debug}[1]{}
%\newcommand{\debug}[1]{{\tiny{(\pageref{#1})}}\normalsize}

\newcommand{\crefSubFigRef}[2]{\namecref{#1}~\hyperref[#1]{\ref{#1}#2}}
\newcommand{\ccite}[1]{Ref.~\cite{#1}} % Ref. [1] instead of [1]
\newcommand{\Py}{\xref{permalloy}}
%%vector
% \newcommand{\vc}[1]{\textbf{#1}}
\newcommand{\vc}[1]{\boldsymbol{\mathbf{#1}}} % Use this instead of \vb, which is deprecated anyway alongside the "physics" package
\newcommand{\circled}[1]{\raisebox{.5pt}{\textcircled{\raisebox{-.9pt} {#1}}}}
\newcommand{\sub}[1]{\ensuremath{_{\mathrm{#1}}}}
\newcommand{\super}[1]{\ensuremath{^{\mathrm{#1}}}}
\newcommand{\av}[1]{\langle #1 \rangle}
\newcommand{\sav}[1]{\langle\hspace{-0.4mm}\langle #1 \rangle\hspace{-0.4mm}\rangle}
%\newcommand{\sav}[1]{\left<\hspace{-0.4mm}\left< #1 \right>\hspace{-0.4mm}\right>}
\newcommand{\tpar}[1]{\frac{\partial #1}{\partial t}}%partial derivative to time
\newcommand{\tder}[1]{\frac{\mathrm{d} #1}{\mathrm{dt}}}%full derivative to time

\newcommand{\EB}{E_\mathrm{B}}
\newcommand{\EBeff}{\widetilde{\EB}}
\newcommand{\EBeffLeft}{\EBeff{}_{,\circlearrowleft}}
\newcommand{\EBeffRight}{\EBeff{}_{,\circlearrowright}}
\newcommand{\EEA}{E_\mathrm{EA}}
\newcommand{\EMC}{E_\mathrm{MC}}
\newcommand{\kB}{\ifmmode k_\mathrm{B} \else $k_\mathrm{B}$ \fi}
\newcommand{\kBT}{\ifmmode k_\mathrm{B}T \else $k_\mathrm{B}T$ \fi}
\newcommand{\clockwise}{\boldsymbol{\circlearrowright}}
\newcommand{\counterclockwise}{\boldsymbol{\circlearrowleft}}
\newcommand{\rmin}{r_\mathrm{min}}
\newcommand{\bigmu}{\makebox{\Large\ensuremath{\upmu}}}
\newcommand{\mavg}{m_\mathrm{avg}}
\newcommand{\qNN}[1][]{q_\mathrm{#1NN}} % Usage: \qNN, \qNN[2], ...

% \newcommand{\abs}[1]{\ensuremath{\left| #1 \right|}\xspace}
\newcommand{\bigsum}{\mathop{\scalebox{2}{\(\sum\)}}}
\newcommand{\bigprod}{\mathop{\scalebox{2}{\(\prod\)}}}
\newcommand{\diff}[2]{\frac{\partial #1}{\partial #2}}			    %% derivative
\newcommand{\dispdiff}[2]{{\partial #1}/{\partial #2}}			    %% derivative
\newcommand{\difff}[2]{\frac{\partial^2 #1}{\partial #2^2}}			%% second derivative
\newcommand{\ddt}[1]{\ensuremath{\partial #1 / \partial t}}
%\newcommand{\d}{\ensuremath{\mathrm{d}}}
\newcommand{\ddiff}[2]{\frac{\d #1}{\d #2}}
\newcommand{\lapl}{\ensuremath{{\Delta}}}
\newcommand{\eqdef}{\stackrel{\small{\mathrm{def}}}{=}}
\newcommand{\inprod}{\cdot}
\newcommand{\bcdot}{\boldsymbol{\cdot}}
\renewcommand{\quote}[1]{``#1''}
\newcommand{\etal}{\textit{et al.}\xspace}

\newcommand{\code}[1]{\textsf{#1}}
\newcommand{\shell}[1]{{\texttt{#1}}\smallskip\\}
\newcommand{\python}[1]{\lstinline{#1}}

\newcommand{\hide}[1]{}

\newcommand{\glijbaantje}[2]{\vspace{-0.5cm}\begin{flushright}{\textit{#1}\\--- #2\\}\vspace{0.5cm}\end{flushright}}
\newcommand{\publicatie}[1]{Material from this chapter has also been published in \cite{#1}.}

\newcommand{\mumax}{mumax\textsuperscript{3}\xspace}
\newcommand{\mumaxplus}{mumax\textsuperscript{+}\xspace}
\newcommand{\oommf}{\textsc{oommf}\xspace}
\newcommand{\nmag}{Nmag\xspace}
\newcommand{\hotspice}{\textsc{Hotspice}\xspace}
\newcommand{\spinengine}{SpinENGINE\xspace}

%%%% EXTRA SYMBOLS
\DeclareSymbolFont{extraitalic}      {U}{zavm}{m}{it}
\DeclareMathSymbol{\Qoppa}{\mathord}{extraitalic}{161}
\DeclareMathSymbol{\qoppa}{\mathord}{extraitalic}{162}
\DeclareMathSymbol{\Sampi}{\mathord}{extraitalic}{165}
\DeclareMathSymbol{\sampi}{\mathord}{extraitalic}{166}
\DeclareMathSymbol{\Stigma}{\mathord}{extraitalic}{167}
\DeclareMathSymbol{\stigma}{\mathord}{extraitalic}{168}

\newcommand{\eqn}[1]{equation \ref{#1}\xspace}
\newcommand{\secateur}[1]{section \ref{#1} (page \pageref{#1})\xspace}
\newcommand{\figuur}[1]{figure \ref{#1} (page \pageref{#1})\xspace}
\newcommand{\Figuur}[1]{Figure \ref{#1} (page \pageref{#1})\xspace}


\newcommand{\kolommen}[2]{
	\begin{minipage}{0.445\textwidth}
		#1
	\end{minipage}
	\begin{minipage}{0.445\textwidth}
		#2
	\end{minipage}
	\smallskip
}

\newcommand{\smallfig}[3]{
	\smallskip
	\begin{center}
		\includegraphics[width=#1\linewidth]{#2}\footnotesize{\\\textsf{(#3)}}
	\end{center}
	\smallskip
}

% 2 args: caption, tabular environment
% 1 optional arg: label
\newcommand{\xtable}[3][nolabel]{
	\jonathanskip\vspace{1em}
	\begin{tikzpicture}
		\node[rounded corners=10pt, fill=figcolor, inner sep=10pt] (figbox) {
			\begin{minipage}{0.955\linewidth}
		  		\centering
		  		\textsf{\tabcaption{#2\label{#1}}}
		  		{\sf
					#3
				}
			\end{minipage}
		};
	\end{tikzpicture}
	\jonathanskip\vspace{0.05cm}
}

% 2 args: filename, caption
% 1 optional arg: width
\newcommand{\xfig}[3][1]{
	\jonathanskip\vspace{1em}
	\begin{figure}[H] % Use the figure environment for compatibility
		\centering
		\begin{tikzpicture}
			\node[rounded corners=10pt, fill=figcolor, inner sep=10pt] (figbox) {
				\begin{minipage}{0.955\linewidth}
					\centering
					\includegraphics[width=#1\textwidth]{#2}
					\textsf{\caption{#3\label{#2}}}
				\end{minipage}
			};
		\end{tikzpicture}
	\end{figure}
	\jonathanskip\vspace{0.05cm}
}

% 5 args: filename1, caption1, filename2, caption2, overallcaption
% 1 optional arg: width of subfigures
\newcommand{\xfigs}[6][0.45]{
	\jonathanskip\vspace{1em}
	\begin{figure}[H] % Use the figure environment for compatibility
		\centering
		\begin{tikzpicture}
			\node[rounded corners=10pt, fill=figcolor, inner sep=10pt] (figbox) {
				\begin{minipage}{0.955\linewidth}
					\centering
					\begin{tabular}{cc}
						\small\parbox[c]{0.47\linewidth}{\centering\textsf{\textbf{(a)} #3}} & 
						\small\parbox[c]{0.47\linewidth}{\centering\textsf{\textbf{(b)} #5}}\normalsize \\
						\includegraphics[width=#1\linewidth]{#2} & 
						\includegraphics[width=#1\linewidth]{#4} \\
					\end{tabular}
					\textsf{\caption{#6\label{#2}}} % Ensure compatibility with hypcap
				\end{minipage}
			};
		\end{tikzpicture}
	\end{figure}
	\jonathanskip\vspace{0.05cm}
}

% 3 args: filename1, filename2, caption
% 1 optional arg: width of subfigures
\newcommand{\xfigsnocap}[4][0.45]{
	\jonathanskip\vspace{1em}
	\begin{figure}[H] % Use the figure environment for compatibility
		\centering
		\begin{tikzpicture}
			\node[rounded corners=10pt, fill=figcolor, inner sep=10pt] (figbox) {
				\begin{minipage}{0.955\linewidth}
					\centering
					\begin{tabular}{cc}
						\includegraphics[width=#1\linewidth]{#2} & 
						\includegraphics[width=#1\linewidth]{#3} \\
					\end{tabular}
					\textsf{\caption{#4\label{#2}}} % Ensure compatibility with hypcap
				\end{minipage}
			};
		\end{tikzpicture}
	\end{figure}
	\jonathanskip\vspace{0.05cm}
}

% 2 args: filename, caption
% 1 optional arg: width of subfigures
\newcommand{\sidefig}[3][1]{
	\jonathanskip\vspace{1em}
	\begin{figure}[H] % Use the figure environment for compatibility
		\centering
		\begin{tikzpicture}
			\node[rounded corners=10pt, fill=figcolor, inner sep=10pt] (figbox) {
				\begin{minipage}[c]{0.5\textwidth}
					\centering
					\includegraphics[width=#1\textwidth]{#2}
				\end{minipage}
				\begin{minipage}[c]{0.445\linewidth}
					\centering
					\textsf{\caption{\justifying #3\label{#2}}} % Ensure compatibility with hypcap
				\end{minipage}
			};
		\end{tikzpicture}
	\end{figure}
	\jonathanskip\vspace{0.05cm}
}


%%%% FIGURES BELOW HERE WERE NOT USED NOR ADJUSTED TO MODERN FORM (YET?)
\newcommand{\sidefigbig}[3][1]{
	\jonathanskip
	\colorbox{figcolor}{
	\begin{minipage}[c]{0.65\textwidth}
  		\centering
  		\includegraphics[width=#1\textwidth]{#2}
	\end{minipage}
	\begin{minipage}[c]{0.325\linewidth}
	 \textsf{\figcaption{#3\label{#2}}}
	\end{minipage}
	}
	\jonathanskip\vspace{0.05cm}
}

\newcommand{\makeshiftfig}[2]{
	\jonathanskip
	\colorbox{figcolor}{
		\begin{minipage}{\linewidth}
			\centering
			
			#1
			
			\textsf{\figcaption{#2}}
			
		\end{minipage}
	}
	\jonathanskip
}

\newcommand{\sidefigr}[3][1]{
	\jonathanskip
	\colorbox{figcolor}{
	\begin{minipage}[c]{0.455\linewidth}
	 \textsf{\figcaption{#3\label{#2}}}
	\end{minipage}
	\begin{minipage}[c]{0.5\textwidth}
  		\centering
  		\includegraphics[width=#1\textwidth]{#2}
	\end{minipage}
	}
	\jonathanskip\vspace{0.05cm}
}

\newcommand{\xfigssmallbig}[6][unused]{
	\jonathanskip
\colorbox{figcolor}{
	\begin{minipage}{0.955\linewidth}
  		\centering
		\begin{tabular}{cc}
		\small\parbox[c]{0.32\linewidth}{\textsf{\textbf{(a)} #3}} & \small\parbox[c]{0.62\linewidth}{\textsf{\textbf{(b)} #5}}\normalsize\\
		\includegraphics[width=0.3\linewidth]{#2} & \includegraphics[width=0.6\linewidth]{#4}\\
		\end{tabular}
		\textsf{\figcaption{#6\label{#2}}}
  	
	\end{minipage}
	
}\jonathanskip\vspace{0.05cm}}

\newcommand{\xfigssize}[8][unused]{
	\jonathanskip
\colorbox{figcolor}{
	\begin{minipage}{0.955\linewidth}
  		\centering
		\begin{tabular}{cc}
		\small\parbox[c]{#7\linewidth}{\textsf{\textbf{(a)} #3}} & \small\parbox[c]{#8\linewidth}{\textsf{\textbf{(b)} #5}}\normalsize\\
		\includegraphics[width=#7\linewidth]{#2} & \includegraphics[width=#8\linewidth]{#4}\\
		\end{tabular}
		\textsf{\figcaption{#6\label{#2}}}
  	
	\end{minipage}
	
}\jonathanskip\vspace{0.05cm}}

\newcommand{\xfigss}[8][0.45\textwidth]{
	%\\[\intextsep]
	\jonathanskip
	\colorbox{figcolor}{
	\begin{minipage}{\linewidth}
  		\centering
		\begin{tabular}{ccc}
		\small\parbox{0.3\linewidth}{\textsf{\textbf{(a)} #3}} & \parbox{0.3\linewidth}{\textsf{\textbf{(b)} #5}} & \parbox{0.3\linewidth}{\textsf{\textbf{(c)} #7}}\normalsize \\
		\includegraphics[width=0.3\linewidth]{#2} & \includegraphics[width=0.3\linewidth]{#4}& \includegraphics[width=0.3\linewidth]{#6}\\
		\end{tabular}
		\textsf{\figcaption{#8\label{#2}}}
  	
	\end{minipage}
	}
	\jonathanskip\vspace{0.05cm}
}

\newcommand{\xfigsss}[9]{
	%\\[\intextsep]
	\jonathanskip
\colorbox{figcolor}{
	\begin{minipage}{0.95\linewidth}
  		\centering
		\begin{tabular}{cccc}
		\small\parbox{0.23\linewidth}{\textsf{\textbf{(a)} #2}} & \small\parbox{0.23\linewidth}{\textsf{\textbf{(b)} #4}} & \small\parbox{0.23\linewidth}{\textsf{\textbf{(c)} #6}} & \small\parbox{0.23\linewidth}{\textsf{\textbf{(d)} #8}} \normalsize\\
		\includegraphics[width=0.23\linewidth]{#1} & \includegraphics[width=0.23\linewidth]{#3} & \includegraphics[width=0.23\linewidth]{#5}& \includegraphics[width=0.23\linewidth]{#7}\\
		\end{tabular}
		
		\textsf{\figcaption{#9\label{#1}}}
	\end{minipage}
	}
	\jonathanskip\vspace{0.05cm}
}

\newcommand{\moviefour}[5]{
	%\\[\intextsep]
	\jonathanskip
\colorbox{figcolor}{
	\begin{minipage}{0.95\linewidth}
  		\centering
			\includegraphics[width=0.21\linewidth]{#1} \hspace{0.1cm} \includegraphics[width=0.21\linewidth]{#2} \hspace{0.1cm} \includegraphics[width=0.21\linewidth]{#3} \hspace{0.1cm} \includegraphics[width=0.21\linewidth]{#4} \hspace{0.1cm}
		%\vspace{-1cm}
		\textsf{\figcaption{#5\label{#1}}}
	\end{minipage}
	}
	\jonathanskip\vspace{0.05cm}
}


\newcommand{\gridfour}[6]{
	%\\[\intextsep]
	\jonathanskip
\colorbox{figcolor}{
	\begin{minipage}{#6\linewidth}
  		\centering
			\includegraphics[width=0.46\linewidth]{#1} \hspace{0.1cm} \includegraphics[width=0.46\linewidth]{#2} \hspace{0.1cm}\\ \includegraphics[width=0.46\linewidth]{#3} \hspace{0.1cm} \includegraphics[width=0.46\linewidth]{#4} \hspace{0.1cm} 
		
		\textsf{\figcaption{#5\label{#1}}}
	\end{minipage}
	}
	\jonathanskip\vspace{0.05cm}
}

\newcommand{\cols}[2]{\parbox{0.5\linewidth}{#1}\parbox{0.5\linewidth}{#2}}

\newcommand{\phdtitle}{Modelling of Reservoir Computing in Perpendicular-Anisotropy Artificial Spin Ice}
\newcommand{\phdauthor}{Jonathan Maes}

%%%% PACKAGES SETUP COMMANDS
\graphicspath{{fig/}} % Package "graphicx"
\sisetup{ % Package "siunitx"
   detect-weight=true
}
\usephysicsmodule{ab,ab.legacy} % Package "physics2": needs to load modules explicitly to avoid namespace conflicts
\hypersetup{ % Package "hyperref": configures options for hyperlinks, bookmarks, and metadata in PDFs.
   plainpages=false,
   colorlinks=true,
   citecolor=gray,
   linkcolor=ugentblue,
   filecolor=black,
   urlcolor=blue!60!black,
   pdftitle={\phdtitle},
   pdfauthor={\phdauthor}
}
\urlstyle{same}
% \captionsetup{ % Package "caption": adjusts the appearance and formatting of captions.
%    width=.85\columnwidth,
%    %format=hang, % 'Figure n' in front of whole caption
%    labelsep=endash,
%    justification=centerlast,
%    labelfont={bf}
% }
\mdfsetup{ % Package "mdframed": margins, padding, and spacing for custom-framed environments.
   skipabove=\topskip,
   skipbelow=\topskip
}
\setstcolor{red} % Package "soul": set color for text highlighting or strikethrough effects.
\titleformat{\paragraph}[hang]{\normalfont\normalsize\bfseries}{\theparagraph}{1em}{}
\titlespacing*{\paragraph}{0pt}{3.25ex plus 1ex minus .2ex}{0.5em}
\newcolumntype{P}[1]{>{\RaggedRight\arraybackslash}p{#1}} % Package "array" & "ragged2e"
\setcounter{secnumdepth}{3} % For "titlesec" package (secnumdepth): adjusts the depth for numbered sections and subsections.
\setcounter{tocdepth}{1} % No excessive table of contents!
\raggedbottom


\setlength{\headheight}{13.6pt} 


%%%% DEFINE BOXED ENVIRONMENTS WITH MDFRAMED PACKAGE
%% Algorithm
\mdfdefinestyle{algorithmstyle}{%
	roundcorner=10pt,
	backgroundcolor=figcolor,
	frametitlebackgroundcolor=gray!30,
	linecolor=gray!50, linewidth=0pt,
	frametitlerule=false,%
	innertopmargin=\topskip,
	innerbottommargin=\topskip,
	subtitleaboveline=true,
	subtitlebackgroundcolor=gray!10,
	%topline=false,
	%bottomline=false,
}
\mdtheorem[style=algorithmstyle]{algorithm}{Algorithm}

%%%% SETUP HEADER
\usepackage{fancyhdr}
\pagestyle{fancy}
\fancyhf{}
\fancyhf[HL]{\nouppercase{{\textsf{\jonathan}}}}  % Links in header het hoofdstuk,
\fancyhf[HR]{\textsf\textbf{{\thepage}}} 						% Rechts het paginanummer.


%%%% CODE HIGHLIGHTING
\definecolor{Code}{rgb}{0,0,0} 
\definecolor{Decorators}{rgb}{0.5,0.5,0.5} 
\definecolor{Numbers}{rgb}{0.5,0,0} 
\definecolor{MatchingBrackets}{rgb}{0.25,0.5,0.5} 
\definecolor{Keywords}{rgb}{0,0,1} 
\definecolor{self}{rgb}{0,0,0} 
\definecolor{Strings}{rgb}{0,0.63,0} 
\definecolor{Comments}{rgb}{0,0.63,0} 
\definecolor{Backquotes}{rgb}{0,0,0} 
\definecolor{Classname}{rgb}{0,0,0} 
\definecolor{FunctionName}{rgb}{0,0,0} 
\definecolor{Operators}{rgb}{0,0,0} 
\definecolor{Background}{rgb}{0.98,0.98,0.98} 
\lstset{
   language=Python,
   breakatwhitespace=false,
   breaklines=true,
   numbers=left, 
   numberstyle=\footnotesize\color{gray}, 
   numbersep=1em, 
   xleftmargin=1em, 
   framextopmargin=2em, 
   framexbottommargin=2em, 
   showspaces=false, 
   showtabs=false, 
   showstringspaces=false, 
   frame=l, 
   tabsize=4,
   basicstyle=\ttfamily\small\setstretch{1}, 
   backgroundcolor=\color{Background},
   commentstyle=\color{Comments}\slshape,
   stringstyle=\color{Strings},
   keywordstyle={\color{Keywords}\bfseries}, 
   keywordstyle={[2]\color{Decorators}\slshape}, 
   emph={self}, 
   emphstyle={\color{self}\slshape}
}

\makeindex
\hyphenation{sen-tence}


%% for b5
%\makeatletter
%\setlength{\oddsidemargin}{7.029mm}
%\setlength{\evensidemargin}{7.029mm}
%\addtolength{\evensidemargin}{21.5mm}
%\setlength{\textheight}{20cm}
%\setlength{\abovecaptionskip}{4pt}
%\makeatother



\begin{document}
\hypersetup{breaklinks=true}
\setlength{\parindent}{0cm}

\frontmatter
\newcommand{\jonathan}{\leftmark}

%\includepdf{fig/coverfig/frontpage}

\thispagestyle{empty}


%\includegraphics[width=22cm]{fig/coverfig/frontpage}\\
 
\fontsize{12pt}{14pt}\selectfont
 
\vspace*{-1.5cm}

\noindent\makebox[\textwidth]{\hfill\includegraphics[width=11cm]{cover/ugentbalk_2022.png}\hfill}

\begin{center}
    \vspace{0.5cm}

    \fontsize{28pt}{28pt}\selectfont
    %\fontsize{17.28pt}{21pt}\selectfont
    \textcolor{ugentblue}{\textsc{Modelling of\\reservoir computing in\\perpendicular-anisotropy\\artificial spin ice}}\\
\end{center}
\vspace{0.0cm}
\fontseries{m}
\fontsize{12pt}{14pt}\selectfont
\vfill

\vspace*{1cm}
\begin{center}
	\centerline{
		\hspace{-.7em}
		\begin{tikzpicture}
			\node[anchor=south west,inner sep=0] (image) at (0,0) {\includegraphics[width=0.95\paperwidth]{cover/OOP_Relaxation.pdf}};
			\fadenode[70pt]{image}
		\end{tikzpicture}
	}
\end{center}
%\newline\null

\vspace*{0cm}
\fontsize{17.28pt}{21pt}\selectfont
\null\hfill \textbf{\textcolor{ugentblue}{\phdauthor}}\hspace*{3.0cm}\\
 
%\vspace{0.5cm}
\fontsize{12pt}{14pt}\selectfont

\newpage

\thispagestyle{empty}
%\noindent
%\includegraphics[width=22cm]{fig/coverfig/frontpage}\\

\fontsize{12pt}{14pt}\selectfont

\vspace*{-1.5cm}
%\includegraphics[width=\textwidth]{ugentbalk} \hspace{1cm} %\includegraphics[height=3cm]{mpilogo}
%\noindent\makebox[\textwidth]{\hfill\includegraphics[width=14cm]{cobined.png} }
\noindent\makebox[\textwidth]{\hfill\includegraphics[width=12.5cm]{cover/ugentbalk_2022.png}\hfill}
\vspace{0.0cm}
\fontsize{20pt}{20pt}\selectfont
\begin{center}
{\textsc{\phdtitle}}\\
\end{center}

\fontsize{16pt}{16pt}\selectfont

\null\hfill Doctoral thesis\\
\null\hfill by \phdauthor\\
\vfill
\fontseries{m}
\fontsize{12pt}{14pt}\selectfont
Thesis submitted to obtain the degree of\\
Doctor of Science: Physics\\
\newline
Public defense: 4th of September, 2025 at\\
Department of Solid State Sciences\\
Faculty of Sciences\\
Ghent University\\


\fontsize{10pt}{12pt}\selectfont
\textbf{Promotors:}\\
Prof. dr. Bartel Van Waeyenberge (promotor)\hfill Ghent University, Belgium\\
Prof. dr. Jonathan Leliaert (co-promotor)\hfill Ghent University, Belgium\\
\newline
\textbf{Examination Commission:}\\
Prof. dr. ir. M. Boone (chair)\hfill Ghent University, Belgium\\
Prof. dr. J. Ryckebusch\hfill Ghent University, Belgium\\
Prof. dr. E. Folven\hfill Norges Teknisk-Naturvitenskapelige Universitet, Norway\\%Norwegian University of Science and Technology, Norway\\
Dr. J. Joos\hfill Ghent University, Belgium\\
Dr. A. Hrabec\hfill Eidgen\"ossische Technische Hochschule Z\"urich, Switzerland\\%Federal Institute of Technology Zurich, Switzerland\\

\newpage

\newpage
\thispagestyle{empty}
\null
\newpage

\clearpage
\thispagestyle{empty}


\chapter*{Acknowledgements}
I acknowledge thee. % TODO
\clearpage


%%%% TABLE OF CONTENTS (on one page)
\begingroup
\makeatletter
% Redefine the \chapter* header macro to remove vertical space
\def\@makeschapterhead#1{%
%\vspace*{50\p@}% Remove the vertical space
{\parindent \z@ \raggedright
\normalfont
\interlinepenalty\@M
\Huge \bfseries  #1\par\nobreak
\vskip 40\p@
   }}
\makeatother

\fontsize{12pt}{15pt}\selectfont
\tableofcontents
\fontsize{12pt}{14pt}\selectfont
\endgroup

\renewcommand{\baselinestretch}{1.1} 	  % De interlinie afstand wat vergroten.
\small\normalsize                        % Nodig om de baselinestretch goed te krijgen.

\mainmatter

\chapter*{Samenvatting (in Dutch)}\label{sec:Preface_NL}
\markboth{Samenvatting (in Dutch)}{Samenvatting (in Dutch)} % For headers in fncychap
\addcontentsline{toc}{chapter}{\nameref{sec:Preface_NL}}
In onze huidige samenleving zijn computers alomtegenwoordig.
Oorspronkelijk werden zij ingezet voor taken die exact uitvoerbaar waren via voorgeprogrammeerde wiskundige operaties, denk bijvoorbeeld aan het verwerken van gestructureerde data, het ondersteunen van het internet, en het uitvoeren van simulaties.
Sinds de opkomst van machinaal leren is hun scala aan mogelijke toepassingen uitgebreid tot veel complexere taken, met als meest recente voorbeeld de grote taalmodellen die ondertussen wijdverspreid zijn.
%Voor het trainen van moderne artifici\"ele systemen, zoals neurale netwerken, is echter een enorme hoeveelheid data nodig.
Het is echter inherent ineffici\"ent om dergelijke modellen te implementeren op conventionele hardware, aangezien de gebruikte algoritmes door vele abstractielagen worden gescheiden van de onderliggende transistor-gebaseerde hardware.
Hierdoor is er veel interesse ontstaan voor alternatieve rekenconcepten, onder andere ook door de gelimiteerde schaalbaarheid van transistoren en de scheiding tussen geheugen en rekenkracht in de traditionele von Neumann computerarchitectuur.
E\'en zulk concept is neuromorphic computing (neuromorfisch rekenen), waarbij men deze artifici\"ele scheidingen wil aanpakken door de synaptische interconnectiviteit van een biologisch brein na te bootsen, waar geheugen en rekenkracht samenvallen.
Een interessant onderzoeksdomein binnen deze context is in-materia computing, waarbij men zich afvraagt welke berekeningen een fysisch materiaal van nature uitvoert en hoe deze kunnen worden benut.
Op deze manier kunnen veel effici\"entere systemen worden bedacht door de vele abstractielagen in conventionele computers te vermijden. \\ % In plaats van wiskundige berekeningen op te leggen aan een fysisch object, zoals tot nu toe gedaan met schakelingen van transistoren.

Op de kruising tussen neuromorphic computing en in-materia computing bevindt zich het zogeheten reservoir computing (RC).
Hierbij wordt een niet-lineair dynamisch systeem met vervagende geheugenwerking --- het zogeheten ``reservoir'' --- gebruikt als zwarte doos.
Wanneer het reservoir verstoord wordt door een inputsignaal, genereert het een complexe dynamische respons.
Hierdoor kunnen de inputgegevens eenvoudiger worden verwerkt: het volstaat om een enkele lineaire transformatie te trainen op deze multidimensionale respons (zonder de eigenschappen van het reservoir aan te moeten passen), aangezien de niet-lineaire transformatie door het reservoir zelf wordt voorzien. 
Opmerkelijk is dat vele fysische media, gaande van een emmer water tot supergeleidende elektronica, de benodigde eigenschappen (niet-lineariteit, vervagend geheugen en hoge dimensionaliteit) van nature bezitten, waardoor zij direct als reservoirs kunnen gebruikt worden.
Dit zorgt voor een veel kortere trainingstijd en lager energiegebruik dan conventionele neurale netwerken, aangezien een deel van de berekening wordt uitbesteed aan het fysische substraat. \par
Het onderzoek dat in deze thesis wordt voorgesteld maakte deel uit van het Europese ``\spinengine'' project, dat zich specifiek toespitste op het gebruik van ensembles van nanomagneten voor reservoir computing.
Bijzondere aandacht ging naar artificieel spin-ijs (ASI), gezien dit soort systeem een sterk niet-lineaire respons genereert, terwijl zeer weinig energie verloren gaat tijdens het verwerken van informatie.
ASI bestaat uit een geordend rooster van interagerende bistabiele nanomagneten, waarvan de magnetisatierichting kan wisselen tussen de twee stabiele toestanden onder invloed van nabije magneten of thermische fluctuaties.
Uit concurrerende lokale interacties onstaat emergente dynamica, wat toelaat om ASI te gebruiken als fysisch reservoir.
Deze thesis richt zich specifiek op ASI met een magnetisatie uit het vlak: dergelijke systemen zijn extra aantrekkelijk aangezien er effici\"ente input- en uitleesmethoden voor bestaan.
Alle voorgaande concepten worden in meer detail ge\"introduceerd in hoofdstuk~\ref{ch:Introduction}. \\
%Het doel van deze thesis is te onderzoeken hoe verticaal gemagnetiseerd ASI ingezet kan worden voor reservoir computing.
%Aangezien een experimentele verkenning van de enorme ontwerpruimte van deze systemen tijdrovend zou zijn, is er geopteerd om een simulatie­-toolkit te ontwikkelen die deze vragen snel en flexibel kan beantwoorden. \\

Hoofdstuk~\ref{ch:Hotspice} stelt de open-source \hotspice Monte Carlo-simulator voor, die de dynamica van ASI simuleert met als doel om reservoir computing-strategie\"en te evalueren en optimale systeemparameters te bepalen.
Er wordt een Ising-achtig model gebruikt: elke magneet in het ASI wordt voorgesteld als een punt-dipool die kan wisselen tussen twee stabiele toestanden.
Verschillende verbeteringen op dit model worden voorgesteld, waarvan de nauwkeurigheid vervolgens wordt vergeleken. % met betrekking tot de simulatie van het dynamische gedrag van ASI.
Deze richten zich op drie aspecten van de simulatie:
\begin{itemize}[noitemsep,nolistsep] % With enumitem package
	\item het in rekening brengen van de eindige grootte van nanomagneten voor de berekening van onderlinge magnetostatische interacties,
	\item de afschatting van de effectieve energiebarri\`ere die de twee stabiele toestanden van een nanomagneet scheidt,
	\item het gebruikte Monte Carlo spin-flip-algoritme voor het simuleren van systeemdynamica.
\end{itemize}
In de tweede helft van het hoofdstuk worden implementatiedetails besproken, met een focus op de berekening van de magnetostatische interactie tussen alle magneten door middel van op voorhand berekende convolutiematrices.
Vervolgens wordt de effici\"entie van de simulator ge\"evalueerd, waarbij een methode wordt voorgesteld om de Monte Carlo-simulatie te versnellen door meerdere magneten tegelijk van magnetisatierichting te laten wisselen binnen \'e\'en simulatiestap.
De correcte werking van de software wordt daarna geverifieerd aan de hand van enkele analytisch oplosbare systemen.
Het hoofdstuk eindigt tenslotte met een vergelijking tussen simulaties en enkele eerder gepubliceerde experimenten, aan de hand waarvan de accuraatheid van de verschillende modelverbeteringen kan worden geanalyseerd. \\

In hoofdstuk~\ref{ch:Applications} wordt deze \hotspice simulator ingezet om de reservoir computing-prestaties van zowel thermisch actief als bevroren verticaal-gemagnetiseerd ASI te evalueren.
Het gedrag van deze systemen wordt grotendeels bepaald door de wisselwerking tussen verschillende energiebijdragen, zoals de loodrechte magnetische anisotropie en de vorm-anisotropie van elke nanomagneet.
Door een zorgvuldige keuze van systeemparameters kan thermisch actief ASI verkregen worden, dat spontaan richting de grondtoestand evolueert over een zekere tijdspanne.
Enkele parameters, zoals de magnetostatische koppeling tussen naaste buren en de netto loodrechte anisotropie, hebben een sterke impact op deze evolutie naar de grondtoestand.
Na het effect van deze parameters te analyseren, worden hun waarden bepaald voor een experimenteel vervaardigd ASI door simulaties te vergelijken met het experimenteel geobserveerde relaxatieproces. \par
Hierop voortbouwend wordt dit relaxatieproces benut voor reservoir computing.
De prestaties van dit thermisch active ASI worden beoordeeld op basis van zijn performantie op twee specifieke taken.
Eerst wordt de niet-lineariteit van het reservoir ge\"evalueerd door het een signaaltransformatie te laten uitvoeren.
Vervolgens wordt de geheugencapaciteit getest, door het reservoir de chaotische Mackey-Glass oscillator te laten voorspellen.
Aangezien thermisch actief ASI moeilijk te vervaardigen is, richt het laatste deel van dit hoofdstuk zich op bevroren ASI.
In tegenstelling tot thermisch actieve systemen kan bij bevroren ASI geen globale input worden toegepast, aangezien deze niet spontaan de grondtoestand kan bereiken.
In plaats daarvan wordt een methode ontwikkeld waarbij elke inputcyclus uit twee aparte stappen bestaat, hetgeen gecontroleerde domeinmuurpropagatie mogelijk maakt.
Hierdoor wordt de domeindynamica benut, wat de prestaties van het reservoir verbetert. \\

De belangrijkste resultaten van deze voorbije hoofdstukken worden tenslotte samengevat in hoofdstuk \ref{ch:Conclusion}, waarna deze thesis wordt afgesloten met een algemene beoordeling van de bekomen resultaten en een blik op toekomstig onderzoek.
\cleardoublepage
\chapter*{Preface}
% TODO END: write preface

This thesis is about ... \\

\cref{ch:Introduction} introduces several key concepts ... \\

\cref{ch:Hotspice} presents the open-source \hotspice Monte Carlo simulator for artificial spin ice (ASI), developed with the intent to evaluate reservoir computing strategies and to determine optimal system parameters.
It uses an Ising-like model: each magnet in the ASI is represented by a point dipole that can switch between two stable states.
Several model variants are presented, whose accuracy in simulating the behaviour of ASI is then compared and discussed.
These variants explore three key aspects of the simulation:
\begin{itemize}[noitemsep,nolistsep] % With enumitem package
	\item the choice of Monte Carlo spin-flip algorithm for simulating system dynamics,
	\item accounting for the finite size of magnets when calculating magnetostatic interactions,
	\item estimating the effective energy barrier that separates the stable states of a nanomagnet.
\end{itemize}
In the second half of the chapter, implementation details are discussed, in particular the calculation of the magnetostatic interaction between all magnets, alongside other strategies to improve performance of the Monte Carlo simulation.
The chapter ends by briefly verifying the software, through comparison of simulations against analytical solutions for several solvable systems. \\

In \cref{ch:Applications}, ... \\

In \cref{ch:Conclusion} ...
\cleardoublepage
\chapter{Introduction}\label{ch:Introduction}
\glijbaantje{It is important to draw wisdom from many different places.\\If we take it from only one place, it becomes rigid and stale.}{Uncle Iroh}
%\glijbaantje{If a machine is expected to be infallible then it cannot also be intelligent.}{Alan Turing}
% TODO: write CMOS instead of or alongside "transistor"

This introduction opens with the broader context and motivation for the research presented in this thesis.
We then introduce the physical reservoir computing paradigm and the key metrics used to quantify performance.
Since we focus on magnetic reservoirs,~\cref{sec:1:Nanomagnetism} describes the fundamentals of nanomagnetism and the physical phenomena that enable efficient input and readout.
Next, we turn to magnetic reservoirs, in particular artificial spin ice, before concluding with a brief description of simulation techniques suitable for these systems.

\section{Context}\label{sec:1:Context}
We live in a society where computers are ubiquitous. % and indispesable.
They have been used for a wide variety of applications, like scientific simulations, graphics rendering, processing well-structured data, running the internet\dots\,
Conventional computers --- consisting of transistors arranged into logic gates --- are well suited for such tasks which require exact mathematical operations to be executed.
However, our world provides many challenging tasks that can not be captured by such simple, exact sets of predefined rules.
The archetypical example of this is driving a car, but this extends far into other complex tasks that are often thought of as requiring some degree of ``intelligence'' --- whatever that word may entail exactly. \par
The field of \idx{machine learning} aims to tackle these rather complex tasks using artificial systems, the most prevalent of which nowadays are \xref{artificial neural networks} that take inspiration from the synaptic connections in a biological brain.
Over the years, neural networks have advanced significantly, most recently evidenced by the rise of large language models which have already had a broad societal impact~\cite{ImprovingLanguageGPT,GPT-4}.
However, training these systems to the point where they can perform any of the aforementioned tasks at a decent level requires vast amounts of data to fine-tune their ever increasing number of internal parameters --- accompanied by an unsustainably rising energy cost~\cite{QuantumNeuromorphicOpportunities,BLOOM_CarbonFootprint_176Bparam}.
This stands in stark contrast to biological brains, which are capable of learning patterns based on a far more limited set of data, while consuming orders of magnitude less power than an artificial system of equivalent computational capability~\cite{NeuromorphicSpintronics}.
Furthermore, running such neural networks on conventional hardware is encountering limits: to increase performance, the number of trainable parameters is rapidly being increased, but transistor scaling and power density limitations make it hard for hardware to keep up. % TODO REF
Hence, while advances in training algorithms and the ever increasing wealth of available data create ever more potent machine learning systems, the underlying approach is leading to a bottleneck. \par
Several factors that contribute to this bottleneck are often quoted in the context of neural networks or other machine learning algorithms running on conventional hardware:
\begin{itemize} % TODO: more references for these bullet points
	\item Conventional hardware uses transistors to perform logical operations.
	The ever increasing computational demands have driven \xlabel{Moore's law}, which states that the number of transistors on a chip doubles every couple of years.
	However, this scaling law appears untenable as transistors approach nanometre scales.
	Relatedly, \xlabel{Dennard scaling} initially allowed the power density of chips to remain roughly constant during this downscaling, but current leakage has inevitably led to increased power dissipation over the past two decades.
	In an attempt to patch this, multi-core processing has become prevalent, but the high power density still limits the number of cores that can operate simultaneously without overheating. % This results in so-called areas of ``dark silicon'' that are unused at a given instant.
	Thus, new kinds of energy-efficient hardware need to be explored.
	\item Conventional computers are built according to the von Neumann architecture, where computation (CPU) is separated from memory (RAM).
	However, transferring information from memory to the processing unit wastes a lot of time and energy, leading to the infamous \idx{von Neumann bottleneck}~\cite{TaskAdaptivePRC}.
	The separation between memory and processing is not fundamentally required: biological brains in the natural world exhibit no such divide as they use a vast network of interconnected neurons, yet they are nonetheless capable of learning, processing and retaining information.
	The field of \idx{in-memory computing} concerns itself with bridging this divide in artificial systems.
	\item Artificial neural networks, like any other machine learning system implemented on conventional hardware, are abstracted far away from the physical domain~\cite{RC_ASI}.
	Most advancements arise from improving their algorithms, but these have to be run on a particular (e.g., von Neumann) computer architecture made up of circuits that handle information encoded by physical devices like transistors~\cite{RC_SuperconductingElectronics}.
	These many layers of abstraction render this approach inherently inefficient, as it does not directly exploit the physical properties of the constituent materials for computation~\cite{RC_ASI}.
\end{itemize}
These various drawbacks of current implementations have become more noticeable over the past decade, as portrayed in~\cref{fig:1:Ngram}.
This has led to the advent of \idx{neuromorphic computing} --- a field of research at the intersection of biology, computer science and physics --- which takes inspiration from the brain to develop more efficient hardware for processing complex data in parallel.
While artificial neural networks are indeed inspired by the brain, their implementation on conventional hardware is typically not considered part of neuromorphic computing, though approaches to implement neural networks on more efficient unconventional hardware are. % rather than applying algorithms on existing transistor-based hardware.
Note that neuromorphic computing does not seek to replace conventional computing: the latter will always find applications in systems that require absolute precision. % think of physical simulations, the transmission of data, or something as common as rendering an image on a screen.
Neuromorphic computing, on the other hand, does not seek exact solutions; in the words of Alan Turing, ``if a machine is expected to be infallible then it cannot also be intelligent.''
%This paradigm shift opens the door to innovative approaches for tackling computational challenges, particularly in areas like real-time sensing, pattern recognition, and prediction.
%Noise can explore a high dimensional landscape related to cost function by escaping local minima.

\vspace{-1.5em}
\sidefig[0.6]{1_Introduction/Ngrams.pdf}{
	\label{fig:1:Ngram}
	Popularity of various key concepts from the year 2000 onwards.
	The surge in popularity of alternative brain-inspired and physical computing methods coincides with raising concerns related to the usage of conventional hardware.
	Data provided by the \href{https://books.google.com/ngrams/graph?content=Neuromorphic computing,Reservoir computing,In-memory computing,von Neumann bottleneck,Dennard scaling,Material computation&year_start=2000&year_end=2022&corpus=en&smoothing=3&case_insensitive=false}{Google Books Ngram viewer} throughout the English corpus, smoothed over a 3-year period.
	\newline
}
\vspace{-1em}

A related concept is \xlabel{material computation}~\cite{NeglectedPillar}, which goes one step further: its goal is to find ways to exploit the inherent computing power present in physical substrates.
Whereas classical computing uses a top-down approach where mathematics is imposed onto a physical material --- for instance to construct logic gates --- material computation works in the opposite direction by instead asking what computations a physical object performs naturally~\cite{RC_ASI}.
By designing materials with specific behaviours, systems can be created that naturally exhibit the desired computational properties, reducing the levels of abstraction that currently exist between physics and computation. \par
Among the various substrates being explored for material computation, magnetic and spintronic (spin transport electronic~\cite{Spintronics}) hardware are especially appealing~\cite{grollier2020neuromorphic,NeuromorphicSpintronicsProspect,QuantumNeuromorphicOpportunities}.
Single-domain ferromagnetic nanomagnets, for example, have the ability to store information without power, respond non-linearly to stimuli~\cite{NeuromorphicSpintronics} and can dissipate very little energy during information processing~\cite{ThermodynamicLimitsComputation,SpintronicsEnergyEfficientComputing}.
The \link{magnetostatic interaction}{magnetostatic (MS) coupling} within ensembles of nanomagnets provides further advantages, as it can give rise to rich collective behaviour like phase transitions, which the constituent magnets would not exhibit in isolation~\cite{NeuromorphicSpintronicsProspect,RC_ASI}.
A well-established example of such an ensemble is \xref{artificial spin ice} (ASI): an ordered lattice of bistable $\approx \SI{100}{\nano\metre}$-size nanomagnets where the lattice geometry does not allow all local interactions to be simultaneously satisfied, leading to emergent dynamics~\cite{ASIpyrochlores}.
The other aforementioned factors have made them a platform of particular interest for scalable and energy-efficient material computation~\cite{PhD_Stromberg}. \par
This finally brings us to the topic of this thesis: ``\emph{\phdtitle}''.
A more detailed description of ASI follows in~\cref{sec:1:ASI}; first,~\cref{sec:1:RC} introduces the general principles of \xref{reservoir computing}.

\newpage
\section{Reservoir computing}\label{sec:1:RC}
% Introduction to RC can take inspiration from: https://arxiv.org/html/2412.13212v1
At the intersection of neuromorphic computing and material computation lies the paradigm of \idx{reservoir computing}.
To explain the underlying concept, we must first take a step back and have a closer look at the aforementioned artificial neural networks.

\subsection{Origin}
\paragraph{Artificial neural networks}
The most ubiquitous machine learning models currently in use revolve around \idx{artificial neural networks} (ANNs).
These were originally inspired by the brain: an ANN consists of a network of non-linear summing nodes (akin to neurons) connected by directed weighted edges (akin to synapses)~\cite{EvaluatingRestrictedESNs,zurada1992introduction,AIsimulateMemoryContinuity}.
Each node holds a numerical value: during operation, this value is determined based on the weighted sum of all nodes that feed into it.
This allows activations to propagate through the network~\cite{lukovsevivcius2009reservoir}.
To prevent these values from accumulating without limit, they are typically bounded to the range $[0,1]$ by applying an \xlabel{activation function} to the weighted sum --- typically a sigmoid, like the logistic map $e^x/(1+e^x)$ or hyperbolic tangent~\cite{SurveyUniversalApproximation}.
A node that uses such an activation function is referred to as an \idx{artificial neuron}.
The non-linearity introduced by the activation function is what enables ANNs to approximate any arbitrary function, allowing them to adequately separate, group and otherwise process input data~\cite{ApproximationFNN,SurveyUniversalApproximation,funahashi1992neural}.
For many real-world tasks, the form of the desired function is unknown, so heuristics are used to ``train'' the weights of all edges based on a given dataset~\cite{EvaluatingRestrictedESNs}.
How exactly this training is done, depends on the structure of the ANN.
Typically, a major distinction is made between \xlabel{feed-forward neural network}s (FNNs) and \xlabel{recurrent neural network}s (RNNs)~\cite{ApproximationRNN}. \par
In a FNN, the nodes are grouped in layers, with connections between all the nodes of two successive layers.
This is schematically illustrated in~\crefSubFigRef{fig:1:NeuralNetworks}{a}: from the blue input node(s) on the left, information propagates to the right through the grey weighted connections, into the green hidden layers, finally yielding one or multiple output values at the red node(s)~\cite{zurada1992introduction}.
FNNs are mainly used for static data processing, such as image classification\footnote{Visual tasks are typically performed using convolutional neural networks --- a particular type of FNN that extracts features from images.}~\cite{ApproximationRNN}.
While~\crefSubFigRef{fig:1:NeuralNetworks}{a} shows a small network of 9 nodes inside 2 hidden layers, many real-world tasks require far larger networks with millions or even billions of weights~\cite{ImageNet,BLOOM_CarbonFootprint_176Bparam,GPT-J-6B}. \par
RNNs are conceptually similar to FNNs, with the important distinction that their connection topology possesses cycles, either in a structured or unstructured manner~\cite{lukovsevivcius2009reservoir,Hopfield1982}.
The existence of cycles has a profound impact, as an RNN may develop self-sustained temporal activation dynamics along its recurrent connection pathways, even in the absence of input.
Hence, mathematically speaking, RNNs are dynamical systems while FNNs are functions~\cite{lukovsevivcius2009reservoir}. % REF verbatim
This enables an RNN to preserve a non-linear transformation of the input history in its internal state, providing dynamical memory and making RNNs especially well-suited for dynamic data processing~\cite{RC_RecentAdvances,ApproximationRNN}.
However, whereas FNNs can be trained using methods such as \xlabel{backpropagation}~\cite{Backpropagation} or gradient descent, the recurrent connections in RNNs make training very computationally expensive and difficult~\cite{DifficultyTrainingRNN,EvaluatingRestrictedESNs,Moon_2021,RC_SuperconductingElectronics,funahashi1992neural}.
Some approaches to train RNNs take inspiration from FNN: back-propagation through time (BPTT), for instance, unfolds the recurrent connections in time to obtain an FNN on which backpropagation can be applied~\cite{jaeger2002tutorial}.
However, this is not so straightforward, leading to issues regarding the calculation of the required gradients and a tendency to get stuck in local minima~\cite{RC_Tensegrity,DifficultyTrainingRNN,D-ESN-Improved}. % Other methods are real-time recurrent learning (RTRL) and extended Kalman filtering (EKE), both explained in jaeger2002tutorial.

\vspace{-1em}
\xfig{1_Introduction/NeuralNetworks.pdf}{
	\label{fig:1:NeuralNetworks}
	Schematic representation of \textbf{(a)} a feed-forward neural network (FNN), \textbf{(b)} a recurrent neural network (RNN) and \textbf{(c)} reservoir computing (RC).
	In the latter, the reservoir is treated as a black box; it can be of mathematical origin, like a RNN, or a physical system like artificial spin ice (ASI).
	Grey connections propagate information from left to right, unless indicated by an arrow.
	Blue (red) nodes represent input (output), while green nodes constitute the hidden layers of the network.
}
\vspace{-1em}

\paragraph{Echo state networks and liquid state machines}
To address the shortcomings of conventional RNN training methods like BPTT, two alternative concepts were independently proposed in the early 2000s, laying the foundation for what would eventually be known as reservoir computing.
Jaeger~\cite{jaeger2001echo} presented echo state networks (ESNs), while Maass~\etal~\cite{maass_LSM} presented liquid state machines (LSMs)\footnote{
	Their names reveal some fundamental desirable properties that the random recurrent network should provide.
	In LSMs, the sparse network of spiking neurons projects inputs into a high-dimensional transient state, akin to waves in a liquid.
	The ``echo'' of ESNs refers to the fading memory provided by the internal dynamics of a recurrent network, where loops can retain information from past inputs.
}.
Both are recurrent networks and operate on very similar principles: their main difference lies in their nodes, as LSMs use \xlabel{spiking neurons} while nodes in ESNs use sigmoidal activation functions.
The crucial property that sets them apart from other training methods is that their underlying recurrent network is randomly initialised and remains unmodified during training~\cite{ReviewESNs,RC_Tensegrity}.
Instead, only a single linear readout layer is trained via linear regression to obtain the desired output from this random network~\cite{D-LSM,D-ESN-Improved}.
This decoupling of internal dynamics from learning is the core insight that underpins both LSMs and ESNs.
Training ESNs and LSMs is a walk in the park as compared to BPTT, as linear regression is far less computationally demanding and does not require the (sometimes problematic) calculation of gradients~\cite{D-ESN-Improved}. \par
From these principles, it is but a small step to reservoir computing (RC), which generalises these ideas by treating the internal network as a black box~\cite{RC_unification,D-ESN-Improved,RC_Tensegrity}.
In the next section we will explore the reservoir computing framework in more detail.

\subsection{Reservoir computing}
Reservoir computing is a machine learning framework suitable for temporal processing tasks~\cite{BookReservoirComputing}.
The concept revolves around a ``reservoir'' --- an unspecified non-linear dynamical system --- that projects the input data into a higher-dimensional space, thereby facilitating the separation of said data by a linear transformation~\cite{appeltant2011information,KUR-24,RC_ASI}. % TODO: use figure like Fig. 2 in ~\cite{appeltant2011information}
The principle is schematically illustrated in \crefSubFigRef{fig:1:NeuralNetworks}{c}.
A weighted sum of the components of (a lower-dimensional representation of) the reservoir's state vector (green nodes), then provides the output (red node) of the system as a whole. % REFs for RC, specifically training a readout layer: Dale et al., 2017; Jensen and Tufte, 2017; Sillin et al., 2013
These weights can be trained with a simple learning algorithm like linear regression in order to produce a desired output~\cite{RC_RecentAdvances, RC_SuperconductingElectronics}.
Crucially, \textit{the properties of the reservoir itself are not modified during training} --- only the final linear readout layer is~\cite{RC_ASI,DynamicEmergence_NanomagneticSystem}. \par
Since the reservoir acts as a black box, its properties can not be changed to obtain the desired response\footnote{
	It is possible to adjust a physical reservoir to exhibit the desired properties for a given task~\cite{AdaptiveProgrammableRC,gartside2022reconfigurable}, but this is different from mathematically tuning the weights of the readout layer via a simple procedure like linear regression.
}.
To solve temporal tasks, a reservoir should be a high-dimensional non-linear system with short-term memory~\cite{NeuromorphicAFMspintronics,RC_RecentAdvances}.
Reservoirs are therefore not limited to being purely mathematical concepts, as many physical systems naturally possess these properties~\cite{RC_DipoleNanomagnets,RC_PassiveFrustratedNM,RC_ASI,RC_RecentAdvances,NeuromorphicOscillators,VowelRecognition4STO,RC_DiffusiveMemristors,RC_MemristorTemporal,gartside2022reconfigurable}, ranging from a bucket of water~\cite{PatternRecognition_Bucket} to photonic systems~\cite{RC_Photonic} and superconducting electronics~\cite{RC_SuperconductingElectronics} --- see~\cref{sec:1:PhysicalRCplatforms}.
This is the concept of \idx{physical reservoir computing} (PRC)~\cite{PRC}, which allows the usage of physical systems without further mathematical abstractions --- apart from the linear readout layer at the end~\cite{RC_RecentAdvances}. \par
While this avoids the many layers of abstraction present in conventional machine learning methods, physical constraints can make it difficult to optimise a physical reservoir~\cite{RC_RecentAdvances}, in contrast to RNNs which have full mathematical freedom.
For instance, when using a physical reservoir to perform a temporal task, the physics of the system are most effectively harnessed when the timescale of system dynamics matches that of incoming data patterns~\cite{KUR-24}, but the timescales of physical phenomena are often hard to control.
Since physical reservoirs can often be hard to scale, multiple techniques have been devised to improve performance by combining multiple reservoirs together~\cite{EvaluatingRestrictedESNs,RotatingNeuronsRC} or by time multiplexing~\cite{appeltant2011information}.
Nonetheless, PRC combines many advantages of all the previously discussed concepts: not much data is required for training, which is computationally very cheap due to the single linear readout layer, enabling low energy usage due to the direct usage of a physical substrate for computation.
The combination of these factors has attracted significant interest to (P)RC.

\paragraph{Reservoir requirements}
For a physical substrate to be suitable as a reservoir, several requirements need to be satisfied: high dimensionality, non-linearity, fading memory, and a separation/approximation property~\cite{appeltant2011information}.
High dimensionality and non-linearity cooperate, as a non-linear mapping of inputs from a small-dimensional into a high-dimensional space can enable the separation of otherwise linearly inseparable inputs, making them both essential properties for classification tasks~\cite{VoltageControlled_SuperparamagneticRC,RC_ASI,appeltant2011information,RC_RecentAdvances}. % The dimensionality is related to the number of independent signals obtained from the reservoir~\cite{RC_RecentAdvances}.
In prediction tasks, the non-linearity of a reservoir can enable the extraction of non-linear dependencies from the inputs for better inference. \par
Fading memory~\cite{boyd1985ApproximatingVolterra} originates from the echo state property of ESNs, where recurrent connections preserve information of the past while it slowly fades out as new input values are applied.
This property ensures that the reservoir state is dependent on recent inputs, while being independent of inputs in the distant past, which is useful behaviour for most temporal processing tasks~\cite{ChaoticTimeSeries_ML,appeltant2011information}. \par % If inputs from the distant past are somehow needed, they can be incorporated using "long short-term memory" (Hochreiter and Schmidhuber).
Finally, the \xlabel{separation property} states that distinct signals should yield distinct responses, while the \xlabel{approximation property} requires similar inputs to remain grouped together in the output space, making the system is insensitive to noise~\cite{RCbenchmarksReview1}.
The former benefits from a chaotic regime, while the latter (and fading memory) are more easily fulfilled in a stable regime~\cite{RC_RecentAdvances}.
As such, many physical reservoirs aim to balance on this ``edge of stability'' whenever it is available in the system, to balance between these properties~\cite{appeltant2011information}.

\paragraph{Mathematical description}
% TODO: add symbols to the figure
In the reservoir computing paradigm, an input signal $s(t)$ perturbs a reservoir $\mathcal{R}$ --- which can be any physical or abstract non-linear dynamical system --- eliciting a response $\mathcal{R}[s](t)$.
The full internal state of the reservoir following this perturbation is often unknown, so a $p$-dimensional readout $\vc{x}(t) \in \mathbb{R}^p$ is assumed (i.e., the green nodes in~\cref{fig:1:NeuralNetworks}).
The final output value\footnote{
	A one-dimensional output is assumed here, but this is not a fundamental restriction: multiple sets of weights can be trained independently if multiple outputs are desired.
} $\hat{y}(t)$ of the system is then obtained by a weighted sum $\hat{y}(t) = \vc{x}^\mathrm{T}(t) \cdot \vc{w}$ of the components of the readout vector $\vc{x}(t)$.
Depending on the task performed, one or multiple such sums can be trained and interpreted in the context of that specific task~\cite{RC_RecentAdvances}. \par % For example, temporal pattern classification may provide multiple outputs, while prediction and generation only uses a single one
When the desired input-output relation $s(t) \mapsto y(t)$ is known for a limited set of data --- the training set --- the weights $\vc{w}$ can be adjusted to minimise the difference between the output $\hat{y}(t)$ and the desired output $y(t)$.
This is usually done by ordinary least-squares (OLS) regression, which requires only a single matrix inversion.
By sampling the readout and desired output at $n$ moments in time, $n$ vectors $\vc{x}_i$ and $\vc{y}_i$, $i = 1, \dots, n$ are obtained.
These can be arranged into an $n \times p$ readout matrix $\vc{X}$ (whose row $i$ is $\vc{x}_i^T$) and the $n$-dimensional desired output vector $\vc{y}$.
The optimal weights can then directly be calculated as
\begin{equation}
	\vc{w} = (\vc{X}^T \vc{X})^{-1} \vc{X}^T \vc{y} \mathrm{,}
\end{equation}
which minimises $\norm{\vc{y} - \hat{\vc{y}}}^2 = \norm{\vc{y} - \vc{X}\vc{w}}^2$~\cite{RC_NNN}.
These optimal weights then remain fixed, and the performance of the reservoir can be evaluated by feeding it the test set --- a previously unknown, yet similar data set as the training set.

\subsection{Metrics} \label{sec:1:RC_metrics}
To compare different physical reservoirs or identify optimal system parameters, metrics are needed that can provide a consistent evaluation of how suitable a given physical system is to be used as a reservoir.
In this respect, two types of metrics are often used: those calculated based on the ranks of matrices, and those calculated through a form of linear regression.
The latter are also referred to as task-agnostic metrics. % Q: find a better name for them?

\subsubsection{Matrix rank-based metrics}\indexlabel{matrix rank-based metrics} \label{sec:1:RC_metrics_KQ}
Legenstein and Maass~\cite{WhatMakesPowerful} proposed three metrics, based on the readout vector $\vc{x}(t)$ of the reservoir.
By periodically sampling this $p$-dimensional vector at $n$ moments in time, an $n \times p$ matrix $\vc{X}$ can be constructed whose rows are these vectors at different moments in time.
Depending on the choice of input sequence, the rank of this matrix gives information about certain properties of the system~\cite{RC_ASI}.

\paragraph{Kernel-quality}
The \idx{kernel-quality} $K$ measures the \xref{separation property}, i.e. how well the reservoir can separate distinct temporal input patterns.
Practically, it is determined by applying $n$ random input sequences and recording the readout after every sequence, each of which consists of a large number of discrete input values (here: 100).
Here, we will use $n=p$ to obtain a square matrix for a more consistent measure.
Using random sequences ensures that they are maximally distinct, such that the rank of this matrix (i.e., the kernel-quality $K$) is a measure for how distinct the responses obtained from the reservoir can be~\cite{Vidamour_2022}.
A more efficient alternative method would be to record the response $\vc{x}$ after each input value, as this would also result in maximally distinct input sequences, though this method strays further from the original definition~\cite{RC_HierarchicalNeuroevolution,RCbenchmarksReview1}.
Hence, a high kernel-quality $K$ indicates that the reservoir exhibits the separation property.
In the extreme case where $K=n$, any desired output can be implemented by the reservoir in combination with its linear readout layer~\cite{WhatMakesPowerful}.

\paragraph{Generalisation-capability}
The \xref{kernel-quality} is insufficient to quantify the computational performance of a reservoir~\cite{WhatMakesPowerful,RC_ASI,IL_Masterproef}.
The \idx{generalisation-capability} $G$ is a similar metric, which measures the \xref{approximation property} i.e., whether similar input sequences yield similar responses.
Practically, it is determined in the same way as the kernel-quality $K$, except that the input sequences are now similar in one way or another.
This can be done by adding a small amount of noise to a given input sequence, but we will follow the method outlined in~\cite{RC_ASI} instead: the first 60 input values will be shared among all sequences, while the remaining 40 are random for each sequence individually.
This way, the rank of the resulting matrix (i.e., the generalisation-capability $G$) measures how distinct the responses of similar inputs are.
A low value of $G$ is therefore desirable, as it indicates robustness to noise and overfitting, as the reservoir adheres to the approximation property~\cite{RCbenchmarksReview1}.

\paragraph{Computing capacity}
To obtain a single metric for the computational capability of a reservoir, Legenstein and Maass proposed to simply subtract the normalised values of both these metrics.
This yields the \idx{computing capacity} $C = K - G$, where higher values typically indicate better performance.
However, care should be taken with such claims, as high computing capacity does not necessarily mean that the reservoir will perform well on all tasks it is given.

\subsubsection{Task-agnostic metrics}
While the matrix rank-based metrics are, strictly speaking, also task-agnostic, the metrics to be discussed now are more commonly referred to when talking about task-agnostic metrics, since there is a clearer relation between them and the types of tasks that the reservoir will perform well on.
A practical definition of non-linearity and memory capacity was given by Love~\etal~\cite{RC_TaskAgnosticMetrics_v2}, which will also be used here.

\paragraph{Non-linearity}
The non-linearity of the reservoir can be measured by considering that its response consists of a linear and a non-linear contribution.
The linear contribution can be extracted by training a linear estimator $\hat{\vc{x}}(t)$ for the current reservoir readout based on the past input $s(t - \tau)$.
More specifically, the estimators
\begin{equation}
	\hat{x}_i(t) = c_i + \sum_{\tau=0}^{k} w_{i,\tau} s(t - \tau) \mathrm{,}
\end{equation}
are trained by adjusting the weights $w_{i,\tau}$ and $c_i$ based on the known discrete-time input-output relationship (on the training set).
The cut-off time $k$ must be longer than the relaxation time of the reservoir --- we will use $k=10$.
By measuring the quality of this estimate using the $\mathrm{R}^2$ correlation coefficient, a measure of the linearity of each separate readout node is obtained.
This can be converted to an overall measure for non-linearity as
\begin{equation}
	\mathrm{NL} = 1 - \frac{1}{p} \sum_{i=1}^p \mathrm{R}^2[\hat{x}_i;x_i] = 1 - \frac{1}{p} \sum_{i=1}^p \frac{\mathrm{cov}^2(\hat{x}_i; x_i)}{\sigma^2(\hat{x}_i) \sigma^2(x_i)} \mathrm{.}
\end{equation}

\paragraph{Linear memory capacity}
The linear memory capacity can also be calculated through by means of linear estimators, but by working the other way around: an estimator for the past input $s(t - \tau)$ is now trained based on the current readout $\vc{x}(t)$ of the reservoir~\cite{tsunegi2019STOforcedsyncRC,NeuromorphicFewShot}.
In other words, the estimator
\begin{equation}
	\hat{s}(t - \tau) = c_\tau + \sum_{i=1}^{p} w_{i,\tau} x_i(t)
\end{equation}
is now trained.
The linear memory capacity is then obtained by summing the $\mathrm{R}^2$ correlation coefficient of all estimators up to a cut-off time $k$:
\begin{equation}
	\label{eq:1:MC}
	\mathrm{LMC} = \sum_{\tau = 1}^{k} \mathrm{R}^2[\hat{s}(t - \tau); s(t - \tau)] = \sum_{\tau = 1}^{k} \frac{\mathrm{cov}^2(\hat{s}_i; r_i)}{\sigma^2(\hat{s}_i) \sigma^2(s_i)}
\end{equation}
Note that the memory capacity can be increased arbitrarily high, by choosing a high cut-off time $k$.
However, the memory capacity typically exhibits a plateau beyond a certain value of $k$, before rising dramatically only when $k$ approaches the size of the training set.
Hence, $k$ should be chosen at this plateau; for the systems we will consider, $k=10$ suffices. \par
Note that this measures short-term memory; long-term memory often refers to the knowledge engrained in the weights of an ANN i.e., the linear readout layer in the case of RC.
Furthermore, this only captures the linear memory capacity: the non-linear memory capacity is far more computationally expensive as it requires a summation over polynomials of different degrees~\cite{RCbenchmarksReview1}.
We therefore limit this discussion to the linear memory capacity.

\paragraph{Parity check}
Closely related to the memory capacity is the parity check metric~\cite{hon2021numerical,tsunegi2019STOforcedsyncRC}.
This is particularly intended for use with binary input, as it defines the parity of a binary input sequence $s(t)$ as
\begin{equation}
	\pi_\tau(t) = \Bigg[\sum_{j = 0}^{\tau} s(t - j) \Bigg] \mathrm{mod}~2 \mathrm{.}
\end{equation}
The parity check (PC) is then defined in the same way as the memory capacity in~\cref{eq:1:MC}, except with $\pi_\tau(t)$ instead of $s(t - \tau)$.
This introduces a form of non-linearity into it, though it still mostly serves as a measure of the memory capacity in binary systems.

\subsubsection{Benchmarks} % Memorisation, frequency generation, classification
In addition to these six metrics, reservoirs are often used to perform specific benchmark tasks in literature.
Example of these include signal transformation and prediction~\cite{AdaptiveProgrammableRC,jaeger2001echo,JaegerHaasWireless,RC_MemristorTemporal,Sunspots_Shougat,gartside2022reconfigurable,appeltant2011information,Vidamour2023}, or speech and image recognition~\cite{farronato2022reservoir,grollier2020neuromorphic,DynamicEmergence_NanomagneticSystem,Vidamour2023}. % See https://arxiv.org/html/2405.06561v
Wringer~\etal~\cite{RCbenchmarksReview1} provide a comprehensive overview of various benchmarks used throughout literature.
They note that --- while less general than metrics --- benchmarks enable a more concrete comparison of reservoirs, as opposed to the more vague notion that two reservoirs are simply different based on their metrics.

\subsection{Physical reservoir platforms}\label{sec:1:PhysicalRCplatforms}
The first true example of physical reservoir computing came in the form of a bucket of water: this proof-of-concept provided input through actuators, and the resulting wave patterns were successfully used for speech recognition~\cite{PatternRecognition_Bucket}.
Similarly, gas bubbles in water have served as an effective reservoir~\cite{RC_GasBubbles}.
While these examples highlight the universal nature of reservoir computing, such macroscopic systems are not very practical.
To allow on-chip integration and achieve low energy consumption, a wide variety of microscopic systems have been exploited as reservoirs~\cite{RC_RecentAdvances}. \par
Traditional machine learning algorithms, like FNN and RNN, are often implemented on traditional transistor-based hardware.
However, for RC specifically, simpler electronic devices exist that can potentially save energy and enable faster computation~\cite{RC_RecentAdvances}.
Memristors are one well-established example: they can store information in their resistance by changing their internal ion distribution, allowing them to act analogously to biological neurons~\cite{MemristorArtificialNeuron,MemristiveNN,RC_DiffusiveMemristors,RC_MemristorTemporal}.
Networks of nanowires are also interesting, as they can self-assemble into a complex recurrent network with memristive junctions forming at intersections~\cite{RC_NNN}.
However, memristors are rather slow.
A faster reservoir can be made of field-programmable gate arrays (FPGAs), which are especially suitable for integer-valued tasks due to their digital nature~\cite{RC_RapidTimeSeries,BookReservoirComputing}. % Autonomous Boolean logic
For instance, FPGAs can efficiently implement cellular automata reservoirs, like elementary rule 90 that exhibits chaotic behaviour and can be implemented solely using XOR gates~\cite{RC_CA}.
Superconducting materials offer low dissipation and extremely fast operation ($\gg \SI{}{\giga\hertz}$), with a number of possible architectures, like for instance a Josephson transmission line~\cite{RC_SuperconductingElectronics}. \par % See also: analogue electronic circuits [TaskAdaptivePRC, Ref. 10], ferroelectrics [TaskAdaptivePRC, Ref. 13], VLSI
Photonic principles have also been exploited to construct reservoirs with very fast operation and low power consumption~\cite{RC_RecentAdvances}.
They make use of on-chip structures like optical waveguides, amplifiers and delay lines~\cite{RC_RecentAdvances}.
One option is to use a single non-linear node, and use delay lines to perform time-multiplexing~\cite{appeltant2011information,RC_AllOptical}.
Another option is to use arrays with non-linear components and delay lines to introduce phase shifts between interconnections, enabling computation through the interference of light~\cite{RC_Photonic,RC_PhotonicSi}. \par
Mechanical systems can also be used as reservoirs, but are not as fast or efficient as the aforementioned systems.
However, mechanical systems provide an extreme example of physical reservoir computing, as it has been shown that the body of a tensegrity robot (tensile structures arranged in trusses, consisting of both springs and incompressible bars) can be used for both locomotion and computation~\cite{RC_Tensegrity}.
Back to the microscopic scale, vibrational modes in micro-electronic mechanical systems (MEMS) can provide the required non-linearity that a reservoir needs.
It is also possible to use chemical processes for computation, for example using processes inspired by those in living cells~\cite{NanoscaleRC,ElectrochemicalPRC,RC_Chemical}. \par
Finally, a wide variety of magnetic and spintronic systems have been proposed for use as physical reservoirs.
Since this thesis focuses on nanomagnetic systems, it is warranted to first explain the physical principles of nanomagnetism before continuing this discussion in~\cref{sec:1_RC_magnetic}.

\newpage
\section{Nanomagnetism} \label{sec:1:Nanomagnetism}
\subsection{Physics of nanomagnets}
This thesis focuses on the usage of magnetic systems as reservoirs.
Generally, magnetism and magnetic fields originate from the movement or acceleration of electric charges.
An object is said to have a \idx{magnetic moment} $\vc{\mu}$ if it experiences a torque $\vc{\tau} = \vc{\mu} \times \vc{B}$ when placed in a magnetic field $\vc{B}$~\cite{IntroMagneticMaterials}.
Conversely, a magnetic moment also generates its own magnetic field
\begin{equation}
	\vc{B}(\vc{r}) = \frac{\mu_0}{4 \pi} \ab[\frac{3 \vc{r} (\vc{\mu} \cdot \vc{r})}{r^5} - \frac{\vc{\mu}}{r^3}] \mathrm{,}
\end{equation}
giving rise to the \xref{magnetostatic interaction} through which magnetic moments influence each other. % (\cref{eq:2:E_MS})
Freely rotatable magnetic moments will therefore prefer a parallel alignment along their connecting axis.

\paragraph{Ferromagnetism}
Some materials, like a bar magnet, exhibit a spontaneous \idx{magnetisation} --- they generate a magnetic field seemingly without any electrical current.
This magnetisation arises at the atomic level, from two distinct sources: the angular momentum of electrons around the nucleus, and the intrinsic spin of the elementary particles comprising an atom~\cite{coey2010magnetism}.
The latter is often dominant, and while the magnetic moments of pairs of electrons often cancel out, atoms with unpaired electrons can exhibit a net magnetic moment~\cite{PhD_Leliaert}. % See slide 35 of https://magnetism.eu/esm/2019/slides/simonet-slides1(magnetism_of_atoms).pdf
However, this does not yet mean that a material composed of such atoms will exhibit a macroscopic magnetic moment like a bar magnet, as the moments of nearby atoms are not necessarily aligned. \par
A non-vanishing net magnetic moment is the privilege of \textit{ferromagnetic}\indexlabel{ferromagnetism} materials.
It is only in these materials that the \idx{Heisenberg-Dirac exchange interaction}~\cite{heisenberg1928theorie} provides a strong tendency for nearby atomic magnetic moments to align, leading to a non-zero macroscopic magnetisation within the material.
This alignment is only possible below a material's \idx{Curie temperature}: at room temperature, \ce{Fe}, \ce{Ni} and \ce{Co} are the only pure elements to be ferromagnetic, though many compounds like \ce{Nd2Fe14B} and \ce{SmCo5} are as well. \par
However, the keyword regarding the exchange interaction is ``nearby'': over longer distances, the magnetisation direction can vary throughout the material as other interactions are also at play --- most importantly, the magnetostatic interaction.
Hence, the magnetisation can be described as a vector field $\vc{M}(\vc{r})$ with a resolution of $\lesssim \SI{5}{\nano\metre}$, as the exchange interaction dominates over shorter length scales.
This forms the basis of micromagnetic theory~\cite{mumax3}: over such small distances, the local magnetisation always has approximately the same magnitude; $\norm{\vc{M}(\vc{r})} \approx M_\mathrm{sat}$, the \idx{saturation magnetisation}.
Applying a magnetic field makes the magnetisation tilt towards the field, but does not increase $M_\mathrm{sat}$: it is an intrinsic property that determines the strongest magnetic moment per volume that a given material can provide.

\paragraph{Single-domain nanomagnets}
Integrating $\vc{M}(\vc{r})$ over the volume of a material yields its net magnetic moment $\vc{\mu}$, but this is almost never equal to $M_\mathrm{sat}V$ since the magnetisation can vary significantly over longer distances.
In fact, an object made of magnetic material often does not exhibit a spontaneous total magnetic moment at all, as the magnetostatic interaction typically causes the magnetisation to loop back on itself to minimise stray fields.
Indeed, in a large piece of ferromagnetic material, it is energetically favourable for \SI{}{\micro\metre}-sized \idx{ferromagnetic domains} to form, whose magnetic moments all cancel each other out when no external field is present~\cite{coey2010magnetism}.
Applying an external field then causes some domains to grow and others to shrink, rather than smoothly canting the magnetisation towards the field.
However, in sufficiently small ferromagnetic particles ($\lesssim \SIrange{100}{500}{\nano\metre}$), the exchange interaction is too strong for separate domains to form~\cite{Kittel_TheoryFMDomains,BrownThermalFluctuations}.
This results in a \idx{single-domain nanomagnet} with a nearly uniform magnetisation $\vc{M}(\vc{r}) \approx \vc{\mu}/V$, whose magnitude is roughly equal to the saturation magnetisation $M_\mathrm{sat}$~\cite{FRENKEL1930,neel1949theorie}. \par % Frenkel and Dorfman first conceived the idea of a single-domain particle based on energy considerations.
When such a nanomagnet exhibits some form of anisotropy, a preferential magnetisation axis will emerge, referred to as the \idx[nolabel]{easy axis}~\cite{nisoli2013colloquium}.
This has led to nanomagnets being used for computation, as the two magnetisation directions along the easy axis can be related to bits `0' and `1'~\cite{MQCA_RoomTemp,NML_Carlton,Gypens_Balanced,Gypens_SelfOrganizing,JM_Masterproef}.
The magnetisation can switch between these two stable states under the influence of nearby nanomagnets, external magnetic fields, thermal fluctuations, electrical currents and a variety of other phenomena~\cite{SwitchingForced_EnergyEfficient,BrownThermalFluctuations,neel1949theorie}.
One common origin of anisotropy is the geometry of a nanomagnet: this \idx[nolabel]{shape anisotropy} is caused by the tendency of magnetic moments to align in the same direction, resulting in a low-energy state when the magnetisation aligns with the longest axis of the geometry~\cite{MagneticCharge}.
Anisotropy can also originate from the crystal structure, but this can be suppressed by an appropriate choice of material; for example, permalloy (\ce{Fe20Ni80}) is an alloy specifically designed to exhibit minimal \xlabel{magnetocrystalline anisotropy}.

\paragraph{In-plane and out-of-plane nanomagnets} % Q: is this OK? What to do with the chapter 3 intro? Just remove 3.1.1?
Two types of nanomagnets are typically distinguished: in-plane (IP) and out-of-plane (OOP) nanomagnets.
IP nanomagnets can readily be obtained by tailoring the geometry of the ferromagnetic particle: a flat geometry guarantees an in-plane magnetisation direction through the \xref{shape anisotropy}.
This easy plane can be reduced to an \xref{easy axis} by using an elongated two-dimensional shape, like an ellipse or rectangle.
OOP magnets are harder to manufacture, as it is often infeasible to create tall structures on wafers.
Therefore, where IP magnets mostly rely on shape anisotropy to create an easy axis, OOP magnets must use other types of anisotropy to favour a magnetisation direction perpendicular to the substrate.
This can be realised using interfacial anisotropy between the ferromagnetic material and the substrate, as for example occurs at the interface between \ce{Co} and \ce{Pt}.
This \idx[nolabel]{perpendicular magnetic anisotropy} (PMA) acts in the immediate vicinity of the interface, where the magnetisation then prefers to align perpendicular to said interface.
The structure of OOP nanomagnets will be discussed in more detail in~\cref{sec:3:OOP_nanomagnet_PMA}~and~\ref{sec:3:E_contributions}.

\subsection{Input and readout for reservoir computing}\label{sec:1:ASI_IO}
To use nanomagnetic systems for RC, methods are required to inject information and read out their resulting magnetic states.
One of the main factors that makes magnetic systems attractive candidates for RC is their ease of electrical interfacing, as many effects can manifest between currents and magnetic materials, enabling efficient input and output.
Electrical control is strongly preferred over external magnetic fields, as it provides greater precision.
In particular, for OOP systems --- the main focus of this thesis --- \xref{spin--orbit torque} provides an efficient input method that synergises well with readout via the \xref{anomalous Hall effect}.

\subsubsection{Input}
\paragraph{Spin-transfer torque}
A well-established method for switching nanomagnets is \idx{spin-transfer torque} (STT)~\cite{SlonczewskiSTT}, as has been widely used in e.g., magnetic random-access memory (STT-MRAM)~\cite{brataas2012current}.
When an electrical current with a spin polarisation $\vc{\sigma}$ enters a ferromagnetic material, the injected spins will exert a torque $\vc{\tau}_\mathrm{STT} \propto \vc{\mu} \times (\vc{\mu} \times \vc{\sigma})$ on the local magnetisation $\vc{\mu}$, causing it to align with the injected spin polarisation.
This polarisation is typically generated by first passing the current through another ferromagnetic layer, whose magnetisation is known or controlled by other means~\cite{SOT_FM_AFM,mumax3tutorial}.
Hence, if both layers are sufficiently close by for the spin polarisation to survive until the second layer, a sufficiently strong STT will cause the magnetisation of the second layer to align with that of the first layer the current passes through.
In OOP nanomagnets, both layers are therefore typically stacked on top of each other, with the current passing vertically through both of them.
While this is an effective method, such a multilayer structure increases fabrication complexity and requires the magnet stack to be fully conductive, which can be limiting in some situations.

\paragraph{Spin--orbit torque}
More recently, switching through \idx{spin--orbit torque} (SOT) has attracted attention as an alternative to STT, as it eliminates the need for a second magnetic layer or vertically flowing current~\cite{mumax3tutorial}.
Instead, it only requires a heavy-metal layer with strong spin--orbit coupling (e.g., \ce{Pt}, \ce{Ta}, \ce{W} \dots) to be adjacent to the ferromagnet~\cite{SOT_Roadmap}.
Typically, a heterostructure is used where a ferromagnet is patterned on top of a heavy-metal underlayer --- for instance, the OOP magnets considered in~\cref{ch:Applications} are \ce{Co}-\ce{Pt} multilayers~\cite{KUR-24}. % exactly the type of structure that supports SOT.
When an in-plane electrical current passes through the underlayer, a spin polarisation will accumulate at the interface with the ferromagnet as a result of the bulk spin Hall effect~\cite{SHE} and/or interfacial Rashba-Edelstein effect~\cite{SOT_Rashba}. % RE effect due to symmetry being broken at an interface. The bulk spin Hall effect often provides the dominant damping-like torque, while any interfacial Rashba effect contributes mainly to the field-like term. % SOT does not originate from a single underlying phenomenon, but is an umbrella term for various effects that all relate to spin--orbit coupling and end up generating a torque in a magnetic material.
These spins, polarised in-plane and perpendicular to the electrical current, can then diffuse into the ferromagnet and exert a torque on its magnetisation~\cite{mumax3tutorial}.
Note that this is indeed distinctly different from STT, where the spin polarisation originates from another magnet and flows straight through the structure~\cite{SOT_Roadmap}. \par
This torque is often decomposed into two orthogonal components.
The field-like torque $\vc{\tau}_\mathrm{FL} \propto \vc{\mu} \times \vc{\sigma}$ causes a precession of the magnetic moment $\vc{\mu}$ around the spin polarisation $\vc{\sigma}$.
On the other hand, the damping-like component $\vc{\tau}_\mathrm{DL} \propto \vc{\mu} \times (\vc{\mu} \times \vc{\sigma})$ aligns the magnetic moment with the spin polarisation~\cite{SOT_Roadmap,SOT_FM_AFM}.
The latter is functionally similar to STT, earning it the name Slonczewski-like SOT.
It has been shown that the field-like SOT can act as a strong transverse effective field, upwards of $\SI{1}{\tesla}$ per $\SI{e8}{\ampere\per\centi\metre\squared}$~\cite{SOT_Rashba}, making it a very efficient method to switch OOP nanomagnets~\cite{vlasov2022optimal,SOTswitchingCoPt,SHE_CurrentInducedSwitching,SpintronicsEnergyEfficientComputing}. \par
However, SOT alone does not guarantee deterministic switching of an OOP ferromagnetic layer, as the in-plane spin polarisation does not differentiate between the `up' and `down' magnetic states.
To resolve this, additional symmetry breaking is required, more specifically with respect to the plane formed by the electrical current and the interface normal~\cite{SOT_Roadmap}.
This can be achieved in a plethora of ways e.g., by an in-plane magnetic field, an exchange bias layer, a source of OOP spin polarisation, exploitation of crystal symmetries if possible, or geometrical asymmetry in the system~\cite{SOT_firstprinciplesCoPt,SOT_Roadmap}.
The polarity of switching depends on the direction of the applied current, the material, and the symmetry breaking.
Hence, when it is known which current direction corresponds to a preferential `up' or `down' magnetisation, electrical control of the magnetisation can be achieved~\cite{SOT_PMAinsulator}.

\subsubsection{Output}
\paragraph{Anomalous Hall effect}
Our method of choice to obtain electrical output from ensembles of out-of-plane nanomagnets is to make use of the \idx{anomalous Hall effect} (AHE).
In the ordinary Hall effect, an electrical current flowing through a conductor (the Hall bar) generates a transverse voltage, since the current generates a magnetic field which exerts a Lorentz force on the electrons.
When the conductor is a ferromagnet, an additional contribution appears that is proportional to the magnetisation of the material.
While its exact origin is still a matter of debate, it is thought to originate from spin--orbit coupling: spin-polarised electrons passing through the ferromagnet get scattered, leading to a charge imbalance across the width of the conductor~\cite{AHE_Culcer,AHE}. \par
Rather than express it as a voltage, the Hall effect is more often described by the transverse resistivity (i.e., the ratio between transverse electric field $E_x$ and longitudinal current density $J_y$) $\rho_{xy} = E_x/J_y = R_0 H_z + R_\mathrm{AHE} M_z$, with the Hall bar in the $xy$-plane~\cite{SHE,AHE}. % TODO: should these symbols be in the symbols list? Especially H vs. B
The first term captures the ordinary Hall effect, the latter the AHE with $M_z$ the OOP magnetisation component and $R_\mathrm{AHE}$ the anomalous Hall coefficient~\cite{AHE}.
Hence, the AHE is greatest when the magnetisation is perpendicular to the conductor plane and has opposite sign for opposite magnetisation directions --- good news for the OOP nanomagnets that we focus on.
This has enabled the usage of AHE for readout in many SOT-MRAM devices, as it can directly be applied in the same heavy-metal/ferromagnet heterostructure as SOT. \par
The readout resolution depends on various factors like $R_\mathrm{AHE}$ and the precise geometry of the Hall bar, with single-magnet resolution achievable. % Typically ~100nm
Since the AHE synergises well with the geometry of OOP nanomagnetic ensembles, it is the method of choice to obtain readout from such reservoirs.

\paragraph{Anisotropic magnetoresistance}
Another phenomenon that can provide readout is \idx{anisotropic magnetoresistance} (AMR): when an electrical current flows through a ferromagnetic material, the resistance weakly depends on the angle $\theta$ between this current and the magnetisation~\cite{AMR}.
More specifically, the total resistance $R(\theta) = R_\perp + (R_\parallel - R_\perp) \cos(2\theta)$~\cite{Hu2023}.
This makes it possible to distinguish between perpendicular and parallel magnetisation states.
While this is a very accessible technique that only requires a measurement of the resistance, it has several limitations. \par
First off, while permalloy exhibits a relatively strong AMR, the absolute effect remains rather small in comparison to the total resistance --- at most a few percent.
This can make it hard to accurately determine the magnetisation state without the use of lock-in amplifiers~\cite{ArchitecturesNanoringRC,Vidamour2023}.
Furthermore, due to the $\cos(2 \theta)$-dependence, AMR does not reveal information about the magnetisation direction, only about the local angle between the magnetisation and current.
It also requires the system to be electrically connected: this makes it an appropriate choice for e.g. the ring system in the \spinengine project, but prevents it from directly being applied to ensembles of nanomagnets.
In interconnected ferromagnets, it only provides a limited amount of information for RC as a lot of the system's dynamics are obscured.
Ensembles of OOP nanomagnets, in particular, would need the current to flow vertically through each magnet, further complicating the system.

\paragraph{Ferromagnetic resonance}
Another possible method to obtain rich output from a reservoir is \idx{ferromagnetic resonance} (FMR)~\cite{AdaptiveProgrammableRC,gartside2022reconfigurable}.
This involves applying a microwave field and measuring the absorption as a function of frequency.
When the frequency matches a resonance of the internal magnetisation dynamics of a nanomagnet, an absorption peak appears.
The complexity of internal magnetisation dynamics, represented in the frequency domain, can provide a rich and non-linear reservoir output~\cite{AdaptiveProgrammableRC,Gomez-Iriarte_FMR}.
However, the peripheral equipment required to perform the technique renders it energy-inefficient and unsuitable for on-chip applications.
Furthermore, as the wavelengths involved in FMR are on the order of centimetres, it does not provide local information --- though for RC purposes it does not need to, since the output space is rich regardless.

\subsubsection{Imaging}
Aside from electrical measurements that provide information about the state of (an ensemble of) magnets, it is also possible to directly image the magnetic configuration of a system through microscopy techniques~\cite{freeman2001advances}.
While these can also be seen as ``output'', they are distinct from the aforementioned electrical measurements as they can not be integrated on-chip.

\paragraph{Magnetic force microscopy}
A relatively common technique is \idx{magnetic force microscopy} (MFM)~\cite{MFM}, which scans a cantilever with a very sharp tip across the sample, similar to atomic force microscopy.
Whereas the latter uses a non-magnetic tip to measure the sample's texture at an atomic scale, MFM uses a tip coated in ferromagnetic material (e.g., \ce{Co}), giving it a vertical magnetic moment.
This enables a measurement of the out-of-plane\footnote{
	It is also possible to measure the in-plane magnetisation by using an in-plane magnetised tip, but this complicates the interpretation as this mixes two components of the magnetisation~\cite{MFM_inplane}.
} magnetic field by scanning the tip at a very low height ($\approx \SI{20}{\nano\metre}$) above the sample~\cite{NML_Carlton,JM_Masterproef}.
However, this does not only measure the magnetic field, but also the atomic forces that a non-magnetised tip would experience.
Therefore, a standard atomic force microscopy scan is first performed to obtain the texture profile of the sample, such that a subsequent MFM measurement can keep the magnetic tip at a constant height above the sample, thereby isolating the magnetic field strength from the atomic forces. \par
The result is a greyscale image, where the brightness indicates the vertical component $B_z$ of the magnetic field close to the sample.
For example, an IP nanomagnet will appear as a pair of regions with bright and dark contrast, because the divergence (convergence) of the magnetic field at its north (south) pole results in opposite vertical magnetic fields on either end~\cite{NML_Carlton}.
An MFM image of OOP magnets (\cref{sec:3:MFM}) is easier to interpret, as there is a one-to-one correspondence between $B_z$ and the the magnetisation state. \par
MFM is relatively cheap to set up as compared to other magnetic microscopy techniques.
It can achieve high resolution, typically on the order of $\lesssim \SI{100}{\nano\metre}$, depending on the height and size of the tip~\cite{MFM}.
However, it is very slow; scanning across a $\SI{}{\micro\metre}$-sized sample can take multiple minutes.
Care must also be taken that the sample is not affected by the magnetic moment of the tip, though this influence is often negligible as the tip is made very small~\cite{Probing_MagnetoOptics}.

\paragraph{Other magnetic microscopy techniques}
While MFM is the most commonly available magnetic microscopy technique, X-ray and optical methods are also often used~\cite{DynamicEmergence_NanomagneticSystem}.
These can provide different information about nanomagnetic systems, as they are directly sensitive to the magnetisation rather than stray fields, making them easier to interpret in systems with in-plane magnetisation.
X-ray based measurements like photo-electron emission microscopy (PEEM)~\cite{PEEM} and magnetic transmission X-ray microscopy (MTXM)~\cite{Imaging_MTXM} can make use of the X-ray magnetic circular dichroism (XMCD) effect to achieve even higher resolution than MFM.
However, they require circularly polarised X-rays, which are typically only available at synchrotron radiation facilities.
Another option is to use optical techniques like the magneto-optical Kerr effect (MOKE)~\cite{KerrFaraday_book}: while these do not require a synchrotron, their resolution is diffraction-limited to several hundred \SI{}{\nano\metre}.
% See also: Polarized neutron reflectometry \cite{DynamicEmergence_NanomagneticSystem}.

\subsection{Magnetic platforms for reservoir computing}\label{sec:1_RC_magnetic}
Besides the physical platforms previously outlined in~\cref{sec:1:PhysicalRCplatforms}, magnetic materials and spintronic phenomena also provide a number of options to construct reservoirs.
They are alluring due to their non-linearity and non-volatility (i.e., memory), as well as low-power electrical interfacing through the various methods outlined in the preceding section.

\subsubsection{Overview of spintronic reservoirs}
Spin-torque nano-oscillators (STOs) are possibly the most abundant devices considered for spintronic reservoirs, owing to their small size and high frequency operation~\cite{tsunegi2019STOforcedsyncRC}.
They are, in essence, a magnetic tunnel junction (MTJ): a stack of two ferromagnetic layers, one fixed and one free, separated by an insulator.
When a current tunnels through the stack, the magnetisation of the free layer will rotate due to STT, at a microwave frequency that depends non-linearly on the input current.
This way, STOs have successfully been used to perform speech recognition~\cite{STO_RC_Riou2021,NeuromorphicOscillators}.
Arrays of STOs are also more broadly used in neuromorphic computing to implement traditional neural networks in hardware~\cite{VowelRecognition4STO}. \par
The magnetisation within ferromagnetic thin films can also exhibit a variety of phenomena, like domain walls, spin waves and skyrmions.
Domain walls~\cite{Venkat_2024} can get stochastically (de)pinned at defects and can mutually interact, which can result in a dynamic equilibrium that is suitable for RC~\cite{DynamicEmergence_NanomagneticSystem}.
They can even exhibit inertia, and the position of a single domain wall can already encode sufficient information for RC~\cite{RC_DW}.
Spin waves excited in ferromagnets with low damping (e.g., YIG) may provide complex interference patterns that can also be used in RC~\cite{RC_SpinWaveInterference}, and can also interact with domain walls.
Skyrmions are non-trivial topological bubble-like features that typically exist at cryogenic temperatures.
Their position and structure can be affected by currents and spin waves, which the skyrmions can in turn affect back, giving rise to a plethora of options to obtain the required non-linearity and fading memory~\cite{RC_TaskAgnosticMetrics_v2,RC_SkyrmionCrystalSW}. \par
This thesis concerns itself with ensembles of magnetostatically coupled nanomagnets.
Due to competing magnetostatic interactions, these ensembles often exhibit collective dynamics like phase transitions and long-range order, which can be exploited for reservoir computing.
A disordered arrangement of nanomagnets has, for example, been used to predict boolean functions like those produced by cellular automata~\cite{RC_PassiveFrustratedNM,RC_DipoleNanomagnets}. % OOP
A class of ensemble of particular interest is artificial spin ice~\cite{RC_ASI}.

\subsubsection{Artificial spin ice}\label{sec:1:ASI}\indexlabel{artificial spin ice}
\textit{Artificial spin ice}\indexlabel{artificial spin ice} (ASI) is a type of metamaterial where magnetostatically coupled single-domain (bistable) nanomagnets are arranged on an ordered lattice in such a way that it is impossible for all local interactions to be satisfied~\cite{RC_ASI,flatspin}.
The magnetisation direction of the constituent nanomagnets can switch between two states under the influence of nearby magnets, externally applied magnetic fields, thermal fluctuations, or any of the electronic input methods previously discussed~\cite{CoerciveFieldReversal,BrownThermalFluctuations,SOT_FM_AFM,brataas2012current}. \par
ASI systems were originally inspired --- and named after --- natural spin ices found in certain crystal structures.
A spin ice is characterised by magnetic \xlabel{frustration}, i.e. the presence of competing interactions which cannot simultaneously all be satisfied.
Spin ices, in turn, got their name from water ice, which is the prototypical system exhibiting geometrical frustration~\cite{nisoli2013colloquium,ZeroPointEntropy,heyderman2013artificial,MagnetizationDynamicsASI}.
Frustration can lead to correlations and collective behavior that the individual elements would not exhibit by themselves, resulting in non-trivial dynamics and complex magnetic orderings~\cite{AdvancesASI,ASI_computation,ApparentFMpinwheel}. \par
ASI is very attractive in comparison to its natural spin ice counterpart.
Modern nano-fabrication methods such as electron beam lithography allow any arbitrary geometry to be realised and many other aspects of the system to be engineered, offering enormous freedom to the designer~\cite{AdvancesASI,ASI_computation}.
Furthermore, the mesoscopic size of ASI (a few hundred nanometers) enables direct observation of their magnetic degrees of freedom through a variety of microscopy techniques and electrical readout methods, in contrast to natural spin ices where imaging individual spins is not feasible~\cite{nisoli2013colloquium,freeman2001advances}.
The aforementioned electrical input methods also enable precise control over these systems.
Hence, the controlled, tunable and easily measurable environment offered by ASI has enabled the study of complex phenomena such as phase transitions with ordered domains~\cite{ApparentFMpinwheel,sklenar2019field,MeltingASI,ImagingBridgedKagome,sendetskyi2019continuous,lou2023competing,branford2012emerging} or glass-like behaviour~\cite{wang2006artificial,ZeroPointEntropy}, vertex-based frustration~\cite{morrison2013unhappy,nisoli2018topologytetris} leading to emergent topological structures such as monopoles and Dirac strings~\cite{ObservationMonopoleASI,mengotti2011kagome}, chiral dynamics~\cite{branford2012emerging,EmergentChiralityRatchet}...
ASI also holds significant potential for computational applications~\cite{heyderman2022spin}, both in conventional logic~\cite{ComputationalLogic_2018,Gypens_Balanced,EngineeringRelaxationComputation} and neuromorphic computing~\cite{ASI_computation,RC_RecentAdvances}. \par
% Q: remove the following paragraph or keep it?
%The study of ASI has since broadened well beyond the scope of frustrated systems, and now the term is used more broadly for ensembles of nanomagnets exhibiting a wide range of physical phenomena~\cite{ASIPathsForward}.
%A vast wealth of both periodic and aperiodic arrangements of nanomagnets has been explored throughout the literature, whose behaviour can vary wildly~\cite{nisoli2013colloquium,heyderman2013artificial,AdvancesASI,ASI_Evolutionary_ALife}.
%For example, a superferromagnetic alignment is most favourable in pinwheel ASI (\crefSubFigRef{fig:2:ASIs}{a,b}), while square-lattice ASI (\crefSubFigRef{fig:2:ASIs}{c,d,i}) prefers an antiferromagnetic (AFM) alignment~\cite{ApparentFMpinwheel}.
%The research presented in this thesis was performed in the context of the European \spinengine project, where three types of magnetic systems were considered: magnetic nanorings~\cite{DynamicEmergence_NanomagneticSystem}, IP ASI~\cite{RC_ASI}, and OOP ASI~\cite{KUR-24}.
%All three present their own advantages and challenges. \par
%Nanoring ensembles show very promising RC performance, but do not provide straightforward on-chip input and readout methods.
%For input, rotating external fields were typically used.
%AMR readout obscures a lot of the system's dynamics and requires the use of lock-in amplifiers, while FMR can provide a higher-dimensional readout in the form of spin-wave spectra~\cite{swindells2024fingerprinting} but requires bulky waveguides. \par % Only provides a single readout value, pretty much
The reservoir computing potential of IP ASI has been demonstrated numerically~\cite{RC_ASI}, but the experimental application of input and state readout is not straightforward.
For OOP ASI, efficient input (SOT) and readout (AMR) mechanisms are available, but their RC potential has not yet been demonstrated.
Therefore, this thesis will fill this gap in knowledge and assess the viability of RC with OOP ASI through numerical simulations.

\newpage
\section{Modelling}\label{sec:1:Modelling}
To determine whether an ASI possesses desirable properties for it to be used as a reservoir, the impact of various system parameters and input methods on reservoir performance must be determined.
Since an experimental exploration of this vast design space would take prohibitively long, a suitable simulation framework is needed.

\subsection{Micromagnetic simulations}
The internal magnetisation dynamics of ferromagnetic structures are the result of the interactions between millions of atomic magnetic moments.
This results in an N-body problem that often has no analytical solution, apart from a select few cases~\cite{PhD_Abert}.
Hence, virtually all practical situations can only be solved in a numerical manner. \par
One particularly useful approximation for the simulation of nanomagnets is \idx{micromagnetic theory}.
It starts from the previously noted observation that the magnetisation within ferromagnets is approximately continuous at length scales below the \xlabel{exchange length} $\lambda \sim \SI{5}{\nano\metre}$, such that it can be represented by a continuous vector field $\vc{M}(\vc{r}) = M_\mathrm{sat} \vc{m}$ with constant norm ($\norm{\vc{m}} = 1$).
Such a continuum approximation is valuable as it enables the simulation of $\SI{}{\micro\metre}$-size systems in a reasonable amount of time through a finite-difference or finite-element discretisation with cells of size $\leq \lambda$.
Examples of finite-difference micromagnetic codes include OOMMF~\cite{OOMMF} and the GPU-accelerated \mumax~\cite{mumax3} and \mumaxplus~\cite{MOR-24}.
Finite-element codes are also available, like Nmag~\cite{Nmag} and magnum.fe~\cite{magnumFE}. \par % Other codes: MagnumNP~\cite{magnumNP} (FD), Boris (FD), tetmag (FE), FeeLLGood (FE)
The main complication in micromagnetic simulations is the infinite range of the magnetostatic interaction, which must therefore be calculated between all cells in the simulation.
This can pose an issue, since the high spatial resolution needed ($\leq \lambda$) can lead to simulations that contain millions of cells.
Furthermore, the time step must be tiny, as magnetisation dynamics can occur over fractions of picoseconds~\cite{PhD_Leliaert}.
Therefore, since ASI may contain hundreds of magnets whose switching rate is often orders of magnitude slower, micromagnetic simulations quickly become infeasible~\cite{leo2021chiral}.
Furthermore, to determine statistically valid RC metrics, many input cycles must be applied --- on the order of 100 for the task-agnostic metrics --- during which many switches occur~\cite{RC_TaskAgnosticMetrics_v2}.
To address these limitations, specialised ASI simulation tools have been developed.

\subsection{Macrospin approximation}\label{sec:1:Modelling_macrospin}
Since the magnetisation of single-domain nanomagnets is often nearly uniform throughout, a \idx{macrospin approximation} can be used where the magnet is represented as a point dipole.
When a nanomagnet has a single easy axis --- as is typically the case in ASI --- it can be approximated as a single Ising spin.
Using such higher-level approximations greatly reduces the computational cost of simulations.
This enables the study of collective behaviour in much larger systems and over far longer timescales, as these simulations are limited by the number of switches rather than the elapsed time.
However, this approach comes at the cost that the internal magnetisation structure of individual nanomagnets is no longer simulated in detail. \par
One example using such a macrospin approximation is the high-performance flatspin simulator~\cite{flatspin}, which implements deterministic spin flipping via a generalised \xlabel{Stoner-Wohlfarth model}~\cite{StonerWohlfarth2008}.
This model revolves around the effective field experienced by each nanomagnet, whose switching is determined by an angle-dependent \xlabel{coercive field} --- it is easier to switch a nanomagnet by applying a field at an oblique angle w.r.t. the easy axis, rather than parallel or perpendicular to it.
These fields are calculated based on a point dipole approximation.
This results in deterministic dynamics, unless temperature effects are included through a random effective thermal field. \par
Besides deterministic simulations, Monte Carlo methods are also often used to simulate both IP~\cite{Qi2008,Cugliandolo2017,LocalizedFrustratedKagome,Brunn2021,Farhan2013,ApparentFMpinwheel} and OOP~\cite{Chioar2014,PerpendicularMagnetizationASI} ASI. % Qi: kagome NN model, Cugliandolo: square vertex-based model, LocalizedFrustratedKagome: kagome NN model, Brunn: IP square 10NN model, ApparnetFMpinwheel: full MS model, Chioar: square magnetostatic, PerpendicularMagnetizationASI: model G and Z.
However, these are typically specialized to a select few lattice geometries.
Early works often accounted only for nearest-neighbour (NN) interactions, whose strength was either arbitrarily set or calculated separately using micromagnetic codes~\cite{Qi2008,PerpendicularMagnetizationASI}. % REFS: 'moller2006artificial': bit broader using dipolar interaction, though seemingly still only for nearest neighbors. 'lou2023competing': dipole model for half-occupation IP Ising. 'sendetskyi2019continuous': dipole model for square ASI. 'EngineeringRelaxationComputation': KMC on dipole model, seemingly for square arrays.
More recently, the importance of long-range magnetostatic interactions has been emphasized to improve correspondence to experiments~\cite{Chioar2014,Rougemaille2011,Brunn2021}.
As such, vertex~\cite{gilbert2014emergent,Saglam2022Tetris,Goryca2021Plasma,MeltingASI} or hybrid NN/vertex charge~\cite{Canals2016,zhang2013crystallites} models have become more prevalent, especially for square-lattice ASI. % 'MeltingASI': 16-vertex ice model, 'gilbert2014emergent': MC sims based on a vertex model and interacting magnetic charges. 'zhang2013crystallites': uses monopoles and NN couplings to model kagome ASI with Metropolis and loop update. 
Beyond-NN interaction models, including full magnetostatically coupled models~\cite{ApparentFMpinwheel,mengotti2011kagome}, are also used for e.g. square-~\cite{Brunn2021,Farhan2013,sklenar2019field} and triangular-based~\cite{Chioar2014,Rougemaille2011,Hofhuis2020} lattices. \par % 'sklenar2019field': 3NN interactions calculated by mumax on quadrupole lattice. 'mengotti2011kagome' use a full dipole model with Ewald summation for PBC.
Here, we opt to develop our own simulation toolkit: \hotspice~\cite{MAES-24}.
This was driven by a desire to blend the approach of flatspin~\cite{flatspin} --- a general-purpose GPU-accelerated IP ASI simulator --- with the rather purpose-built Monte Carlo toolkits listed above, to obtain efficient simulation of RC in OOP ASI.
The result is a versatile Monte Carlo simulator that aims to capture IP and OOP ASI physics with a minimal set of physically motivated parameters (most importantly, the magnets' anisotropy, magnetic moment and temperature) for a variety of lattice configurations.
System dynamics are determined by the switching probability of each magnet, which is calculated based on the energy barrier separating its two stable states.
Given the importance of considering long-range \xref{magnetostatic interactions} --- particularly for OOP systems --- as highlighted by Chioar~\etal~\cite{Chioar2014}, \hotspice{} explicitly accounts for the magnetostatic interactions between all magnets.
Various optional model variants have also been implemented to enable a closer match to experiments.
The details of the \hotspice simulator will be extensively discussed in~\cref{ch:Hotspice}, after which it is used in~\cref{ch:Applications} to explore reservoir computing in OOP ASI.
\cleardoublepage
\chapter{Methods: ``\hotspice'' simulator for ASI}\label{ch:Hotspice}
% \glijbaantje{It's not a bug, it's a feature.}{Someone}

\begin{adjustwidth}{2em}{2em} % TODO END: update once published.
    \begin{center}
        \textbf{Material from this chapter has also been published in:} \\
    \end{center}
    \vspace{1em}
    \begin{adjustwidth}{0em}{2em}
	    \begin{itemize}
	    	\item[\cite{MAES-24}] J.~Maes, D.~De~Gusem, I.~Lateur, J.~Leliaert, A.~Kurenkov, and B.~Van~Waeyenberge.
	    	\newblock The design, verification, and applications of Hotspice: a Monte Carlo simulator for artificial spin ice.
	    	\newblock \emph{ArXiV}, arXiv:\penalty0 2409.05580, 2024.
	    \end{itemize}
    \end{adjustwidth}
    \vspace{1em}
    \begin{center}
        \centering\rule{0.6\linewidth}{0.4pt}
    \end{center}
    \vspace{1em}
\end{adjustwidth}

To assess the potential of \xref{Reservoir Computing} in (perpendicular-anisotropy) \xref{Artificial Spin Ice}, i.e. the central goal of this thesis, simulations will be required to efficiently explore the impact of system parameters on performance. Nanomagnetic systems are often simulated using micromagnetic codes, such as the finite-difference-based \mumax~\cite{mumax3} and \oommf~\cite{OOMMF} or the finite-element-based \nmag~\cite{Nmag}, which capture the magnetization dynamics of individual nanomagnets in great detail. However, the time between successive switches of a nanomagnet is not necessarily similar to the timescale of micromagnetics: when the simulated time extends beyond several microseconds, simulating even a modest number of magnets -- on the order of several dozen -- becomes computationally unfeasible~\cite{leo2021chiral}. \par % Q: is it ok to just copy sentences from my paper?

% TODO: this is placeholder text to test some commands.
Hotspice is written in Python 3.10, and can choose to calculate on either the CPU or GPU.
CuPy v11.4~\cite{CuPy} is used for GPU-accelerated array manipulation, otherwise NumPy and SciPy are used for the CPU.
When writing code for the GPU, it is important to keep the \href{https://docs.nvidia.com/cuda/cuda-c-best-practices-guide/index.html}{Cuda best practices} in mind.
Unfortunately, I only found this page after one year of Hotspice development *shrug*.
Another Python package for parallel programming is the \python{cuNumeric} package by Nvidia, which is also designed to be a drop-in \python{numpy} replacement but supports multiple GPUs, which can be useful for very large simulations that do not fit into the memory of a single GPU.
There is also the \code{nvc++} compiler but this is for Cython code. \par

\section{Model}

\cite{Hotspice}

\xfig[1.0]{2_Hotspice/ASIs.pdf}{\label{fig:2:ASIs}Predefined artificial spin ice (ASI) lattices available in \hotspice{}. The unit cell of each lattice is delineated by a central grey rectangle. The red indicator defines the lattice parameter $a$. In the Ising approximation, the magnetization of in-plane magnets (left) aligns along the major axis of the depicted ellipses. Out-of-plane magnets (right) are illustrated as circles.}\\

\section{Energy calculation}
\subsection{Energy terms}
\subsection{Finite-size correction for magnetostatic energy}
\subsubsection{Second-order correction for dipoles}
\subsubsection{Dumbbell model}
\subsubsection{Comparison}
\subsection{Effective energy barrier}
\subsubsection{Intrinsic barrier from shape anisotropy}
\subsubsection{Mean-barrier approximation} % and relevant choices with discontinuities
\subsubsection{Asymmetric barrier}
\subsubsection{Exact solution} % OOP only

\section{Dynamics}
\subsection{N\'eel relaxation: temporal evolution}
\subsection{Metropolis-Hastings: sampling equilibrium states}
% TODO: decide whether to talk about nomenclature details etc. here, or in the introduction.
\subsection{Other Monte Carlo algorithms} % Not implemented, explain why Wolff could be useful but why it's not in Hotspice

\section{Implementation details}
\subsection{API} % Structure of Hotspice, explain all modules
\subsubsection{ASI}
\subsubsection{RC}
\subsection{Performance} % Several factors impacting performance, explain by using that one graph
\subsubsection{GPU vs. CPU}
\subsubsection{Kernels} % If not yet explained in detail earlier, talk about the perpendicular kernel etc. here
\paragraph{Numerical error with cut-off kernel}
\subsubsection{Multi-switching in Metropolis-Hastings}
\paragraph{Minimal distance between sampled magnets} % Derive equation
\paragraph{Poisson disc}
\paragraph{Grid-select}
\paragraph{Hybrid grid-poisson}
\subsubsection{RNG} % Little to no impact iirc, can be omitted probably
\subsection{Advantages and disadvantages of this approach} % Hindsight is 20/20
\subsubsection{Grid}
\subsubsection{Input/output RC}

\section{Verification} % See paper
\subsection{Hexagon}
\subsection{Non-interacting spin ensemble}
\subsection{Exchange-coupled OOP square system}
\subsubsection{Critical slowing down}
\subsection{Exchange- and magnetostatically-coupled OOP square system}
\subsection{Square-to-pinwheel transition angle}
\cleardoublepage
\chapter{Applications}\label{ch:Applications}
% TODO END: go through this chapter (and the others, probably) at the end and see if any comments with additional information can be put or merged anywhere.
\glijbaantje{If you find that you're spending almost all your time on theory,\\ start turning some attention to practical things; it will improve your theories.\\ If you find that you're spending almost all your time on practice,\\ start turning some attention to theoretical things; it will improve your practice.}{Donald Knuth}

\begin{adjustwidth}{2em}{2em} % TODO END: update once published.
	\begin{center}
		\textbf{Material from this chapter has also been published in:} \\
	\end{center}
	\vspace{1em}
	\begin{adjustwidth}{0em}{1.5em}
		\begin{itemize}
			\item[\cite{KUR-24}] A.~Kurenkov, J.~Maes, A.~Pac, G.~M.~Macauley, B.~Van~Waeyenberge, A.~Hrabec, and L.~J.~Heyderman.
			\newblock Perpendicular-anisotropy artificial spin ice with spontaneous ordering: a platform for neuromorphic computing with flexible timescales.
			\newblock \emph{ArXiV}, arXiv:\penalty0 2408.12182, 2024.
		\end{itemize}
	\end{adjustwidth}
	\vspace{.5em}
	\begin{center}
		\centering\rule{0.7\linewidth}{0.4pt}
	\end{center}
	\vspace{1em}
\end{adjustwidth}

Our main focus will be on out-of-plane systems. First, however, we will rather briefly discuss several in-plane cases.

\section{In-plane systems}
\subsection{Pinwheel reversal}\label{sec:3:IP_Pinwheel_reversal}
\subsection{Kagome reversal} % Q: move these reversals to the Hotspice chapter?
\subsection{RC by clocking in Pinwheel ASI}

\section{Reservoir computing in out-of-plane systems}
The main factor that attracted attention towards OOP ASI for RC, was that they allow very efficient and relatively simple input and read-out methods to be used.
Early on in the \spinengine project, methods to read the state of OOP ASI had already been demonstrated experimentally on a small scale by ETHZ/PSI.
However, their potential for RC had not yet been investigated.
As \hotspice is well-suited for the simulation of OOP ASI, this presented an appropriate use-case for the software developed in~\cref{ch:Hotspice}. % Furthermore, the other simulator developed by the consortium (flatspin) only supports in-plane magnetisation. Furthermore, \hotspice is more accurate for OOP ASI than IP ASI due to the high degree of symmetry in the former.
\par
Therefore, in this section \hotspice will be used to assess the viability and characteristics of a couple of methods to achieve RC in OOP ASI.
It will be shown that RC is indeed possible in such systems by using an appropriate input and readout protocol.
This section starts with a discussion of the basic properties of OOP ASI and the motivation for researching them.
Then, the section is divided into two main parts, discussing non-volatile and thermally active ASI. % At first, the fabricated OOP ASI were not thermally active, though ASI which spontaneously relaxed to a checkerboard state was achieved later.
\subsection{Characteristics of OOP ASI}
\subsubsection{Motivation} % Input and output
While it is possible to use an external field to switch both IP and OOP nanomagnets, this is undesirable in a real device due to stray fields and the need for bulky magnets.
The main allure of OOP ASI stems from the existence of efficient input and readout mechanisms for such OOP magnets.
Input can be applied via spin-orbit torque (SOT)~\cite{SOT_FM_AFM,SOTswitchingCoPt}, and the system can be read out using the anomalous Hall effect (AHE)~\cite{AHE}.
Both of these only require a current to be passed through an underlayer that is electrically connected to the sample, and were explained earlier in the introduction (\cref{sec:1:ASI_IO}).
% Nonetheless, external fields are often used for input anyway in the lab during the first stages of development, because they are easier to handle than the more complicated methods that would be used in a real device. By the end of the project, Alex had used SOT in his actual devices.

\subsubsection{Structure of OOP nanomagnets} % Co/Pt multilayer, energy contributions
\label{sec:3:OOP_nanomagnet_PMA}
Where IP magnets mostly rely on shape anisotropy to create an easy axis, OOP magnets can be realised using interfacial anisotropy between the ferromagnetic material and the substrate.
One example of this is found at the interface between \ce{Co} and \ce{Pt}, as was used in the experimental OOP ASI fabricated within the context of the \spinengine project.
This \idx{perpendicular magnetic anisotropy} (PMA) acts in the immediate vicinity of the interface, where the magnetisation then prefers to align perpendicular to said interface. \par
Typically, a number of \ce{Co}-\ce{Pt} interfaces are stacked vertically, with the number and thickness of \ce{Co} layers then controlling the OOP anisotropy.
The ferromagnetic layers in such stacks are typically made to be very thin for two main reasons.
Firstly, since the PMA only acts at the interface it does not increase for thicker layers.
Secondly, recall that the energy associated with \xref{shape anisotropy} increases with volume ($\EB=K_\mathrm{u}V$), so a maximal number of layers for a given amount of ferromagnetic material is preferable for PMA to counteract shape anisotropy. \par
Other factors contribute to the anisotropy as well, e.g. a non-linear dependence on the \ce{Pt} layer thickness due to the RKKY interaction~\cite{RKKY_RK,RKKY_K,RKKY_Y}.
However, this coupling is likely not dominant in a \ce{Co}-\ce{Pt} stack, as \ce{Pt} nearly satisfies the Stoner criterion~\cite{PtMagneticOrder}.
Since the \ce{Pt} is combined with a ferromagnetic material (\ce{Co}), it is thought that this induces a ferromagnetic coupling between \ce{Co} layers in spite of the RKKY interaction, for reasonable \ce{Pt} thicknesses~\cite{PerpendicularMagnetizationASI}.

\subsubsection{ASI ground state}
The ground state of any system being used for RC is of utmost importance; it dictates which input methods will induce desirable dynamics in the system and which input protocols will fall flat.
The ground state of most OOP spin ices is some form of anti-ferromagnetic (AFM) ordering, because the dipolar interaction will always encourage opposite magnetisation between any two neighbouring OOP magnets.
Depending on the particular lattice upon which the OOP magnets are placed, this may (e.g. triangle/Cairo; \crefSubFigRef{fig:2:ASIs}{j,l}) or may not (e.g. square/honeycomb; \crefSubFigRef{fig:2:ASIs}{i,k}) induce frustration in the system.
If no frustration is present, the symmetry of the system imposes exactly two opposite ground states.
If frustration is present, no real ground state exists apart from local interactions, though the simple two-fold symmetry will still be present, as in any ASI.
For the OOP square-lattice ASI in particular, which is the lattice we will consider throughout this chapter, the ground state is an AFM checkerboard pattern of `up' ($\uparrow$) and `down' ($\downarrow$) magnetisation.

\subsubsection{Input and readout: experiment to simulation}
\paragraph{Input by spin-orbit torque}
In the experimental system, the intent was to switch the OOP magnets using SOT~\cite{SOTswitchingCoPt}.
Compared to an \xref{external magnetic field}, this allows far greater control over the stimulus applied for each input bit, and can more easily be integrated on-chip without generating stray fields. \par
Since \hotspice does not work with torques, but rather with switching energies, SOT can not directly be modelled as a torque.
However, the net desired effect of SOT in our application is deterministic switching of magnets, which can be achieved in the presence of in-plane symmetry breaking~\cite{SOT_Roadmap}. % SOT_Roadmap mentions symmetry breaking several times, very useful resource
Broken symmetry creates a preferential magnetisation direction, which can be modelled in \hotspice as an additional external field $B_z$.
This is a crude approximation to fit SOT into the Ising-like model; the modelled field $B_z$ is not necessarily related to the field-like SOT torque~\cite{SOT_firstprinciplesCoPt}, though the SOT current density is likely non-linearly related to the modelled field $B_z$.
% In the case without symmetry breaking, the effect may be modelled as a temporary reduction of the PMA, as in reality a sufficiently strong SOT (applied over multiple $\tau_0$) will push the magnetic moment in-plane~\cite{SOT_Roadmap}, after which it falls back randomly (with higher likelihood to the lowest energy state, I suppose). But this approach raises multiple additional questions: is it meaningful to reduce the barrier only slightly?

\paragraph{Readout by anomalous Hall effect}
In our paper~\cite{KUR-24}, \textit{A. Kurenkov} experimentally demonstrated electrical readout of the average magnetisation $\llangle s \rrangle$ using the anomalous Hall effect~\cite{AHE}.
This can therefore be used as a magnetic state readout in the simulations.
However, for reservoir computing purposes, multiple readout values are beneficial, so the average magnetisation of each lattice column was used for the readout $y_i(t)$. % Fig. 6a

A discussion about the experimental feasibility of such a grid of local readouts is given in Supplementary Information 7 of our paper~\cite{KUR-24}.
In short, by adding more electrodes to the Hall bar --- the layer underneath the ASI used to measure the magnetic state through the AHE --- inhomogeneities in the current density distribution result in different currents flowing through the different nanomagnets.
This consequently results in unequal contributions to the Hall resistance, providing a way to differentiate between the different magnets.
Individual states of the nanomagnets, not just $\mavg$, can be extracted through a set of linear operations as detailed in the aforementioned Supplementary Information.
In the context of reservoir computing, this means that the use of this readout approach with multiple electrodes can be computationally equivalent to knowing the exact magnetic configuration of the system.


%We will later notice that simple stimuli are not sufficient for RC in this system, and instead more complicated patterns will have to be used to meaningfully address the system. However, the practical implementation of such complicated stimuli requires more complicated current line layouts, which also cannot interfere with the readout. This proves to be a conundrum, which I will address later again.

\subsection{RC in non-volatile OOP square ASI}
\subsubsection{Considerations for devising an appropriate input scheme}
\paragraph{Global field can not work} % Start with schemes that are not appropriate
\paragraph{AFM field can work}
\paragraph{Input encoding} % Binary vs. byte etc. Sigmoid thingy?
\subsubsection{Two-step input scheme: clocking} % With importance of E_B_std

\subsection{Thermally active OOP ASI}
\subsubsection{Energy contributions}
Thermally active OOP ASI are hard to manufacture, because they must strike a delicate balance between several strong energy contributions to end up in an OOP state with a net energy barrier of only several tens of $\kBT$ at most.
To appreciate this difficulty, let us first take a closer look at the various energy contributions in the system --- ranging from local anisotropy to large-scale interactions --- and their dependence on the geometrical structure of the OOP magnets.
\begin{enumerate}
	\item \textbf{Perpendicular magnetic anisotropy} constitutes the largest contribution to the total OOP anisotropy.
	Its interfacial origin and characteristics have been discussed earlier in this chapter (see \cref{sec:3:OOP_nanomagnet_PMA}). % Anisot E: \approx\SI{1.1e-16}{\joule}
	\item \textbf{Shape anisotropy}\indexlabel[nolabel]{shape anisotropy} of each layer due to their \xlabel[nolabel]{demagnetising field}.
	It is of a similar order of magnitude as the PMA, but of opposite sign, favouring an IP magnetisation rather than OOP.
	This contribution was also discussed earlier, see~\cref{sec:2:shape_anisotropy}~and~\ref{sec:3:OOP_nanomagnet_PMA}. % Demag E: \approx\SI{-1.6e-16}{\joule}
	\item \textbf{Inter-layer magnetostatic interaction}.
	The separate FM layers within a single stack also interact, as they can all be considered as separate vertically stacked magnets.
	This vertical stacking provides an additional contribution to the OOP anisotropy because, in general, the magnetisation of separate magnets prefers to point along the axis connecting those magnets.
	One can think of this effect as being similar to the shape anisotropy, but on the scale of separate layers rather than individual atomic magnetic moments.
	This contribution aids the PMA in stabilising the OOP anisotropy, counteracting the \xref{demagnetising field}, but is smaller in magnitude. % Cross MS E: \approx\SI{0.5e-16}{\joule}
	\item \textbf{Inter-magnet magnetostatic interaction} is the only contribution --- besides external input --- that distinguishes between the $\uparrow$ and $\downarrow$ magnetisation states.
	All aforementioned energy contributions are internal to individual OOP magnets, making the \xref{magnetostatic interaction} the sole reason for the AFM ground state of the OOP ASI. % E_MS: \lesssim\SI{1.1e-17}{\joule}
	\item \textbf{Input}.
	In simulations, an \xref{external magnetic field} is used as a proxy to model SOT.
\end{enumerate}
For a system to exhibit spontaneous switching, the first three of these contributions should mostly cancel out --- to the level where their combined influence is of the same order of magnitude as the inter-magnet magnetostatic interaction, which is usually at least an order of magnitude weaker than the other contributions separately.
As we will see in~\cref{sec:3:relaxation}, achieving spontaneous thermal ordering to the ground state may require even finer control.

\paragraph{Magnet geometry}
\label{sec:3:OOP_geometry}
% The geometry of nanomagnets in an OOP ASI is shown schematically in~\cref{fig:3:OOP_geometry}.
% TODO: make own figure of nanomagnets?

The magnets considered here are round with a diameter $D_\mathrm{NM}$ and consist of $N_{\ce{Co}}$ layers of \ce{Co} with a thickness $t_{\ce{Co}}$, separated and surrounded by \ce{Pt} layers of thickness $t_{\ce{Pt}}$.
They are placed in an ASI where the lateral spacing between magnets is $S_\mathrm{ASI}$ --- resulting in a lattice parameter $a=D_\mathrm{NM}+S_\mathrm{ASI}$ in a square lattice. \par
These five geometrical parameters have a profound effect on the balance between the aforementioned energy contributions, as summarised in~\cref{tab:3:interactions_geometry}.
The polynomial dependence of some of these relations makes the system quite hard to control: a slight manufacturing difference can yield a vastly different net OOP anisotropy, possibly inhibiting the spontaneous formation of a checkerboard ground state.
Note that the magnetostatic interactions depend on the square of the ferromagnetic volume $V=n_{\ce{Co}} t_{\ce{Co}} D_\mathrm{NM}^2$.

\xtable[tab:3:interactions_geometry]{
	Dependence of the energy components present in OOP ASI on the geometrical parameters.
	$x$ represents the parameter above the column.
	The magnet centre-to-centre distance $r = D_\mathrm{NM} + S_\mathrm{ASI}$.
	PMA depends non-trivially on $t_{\ce{Pt}}$.
	Empty cells indicate no significant dependence.
}{
	\begin{tabular}{r|c|c|c|c|c}
		\multicolumn{1}{r}{} & \multicolumn{1}{c}{$t_{\ce{Co}}$} & \multicolumn{1}{c}{$n_{\ce{Co}}$} & \multicolumn{1}{c}{$D_\mathrm{NM}$} & \multicolumn{1}{c}{$t_{\ce{Pt}}$} & \multicolumn{1}{c}{$S_\mathrm{ASI}$} \\
		\hline \hline
		PMA (interfacial anisotropy) &  & $x$ & $x^2$ & $\aquarius$ &  \\
		\hline
		Shape anisotropy (demag) & $x^2$ & $x$ & $x^4$ &  &  \\
		\hline
		Inter-layer MS interaction & $x^2$ & $x^2$ & $x^4$ & $x^{(< -3)}$ &  \\
		\hline
		Inter-magnet MS interaction & $x^2$ & $x^2$ & $x^4$ &  & $r^{(< -3)}$ \\ % S_ASI: $>r^{-3} + \frac{3}{16}D_\mathrm{NM}^2 r^{-5} $
		\hline
		Input (external field) & $x$ & $x$ & $x^2$ &  &  \\
		\hline
	\end{tabular}
}

These five geometrical parameters are bounded by several physical and practical constraints.
Note that the numerical values of the limits given below are only indicative and may depend on the exact values of the other geometrical parameters.
\begin{itemize}
	\item \textbf{FM layer thickness} $\boldsymbol{t_{\ce{Co}}} \lesssim \SI{1.45}{\nano\metre}$. \newline
	This upper bound was experimentally determined~\cite{KUR-24} and is imposed by shape anisotropy.
	Layers must be sufficiently thin to maintain OOP anisotropy, as \xref{shape anisotropy} is proportional to volume (and hence $\propto t_{\ce{Co}}$) while the PMA is mostly independent of $t_{\ce{Co}}$. \par % t_{\ce{Co}} controls the demag compensation of $K_u$.
	The closer the thickness is to this upper bound, the lower the effective OOP anisotropy will be.
	Therefore, non-volatile ASI must steer clear of this limit, while thermally active ASI generally try to approach this limit as close as possible.
	As such, this geometrical parameter is often key for balancing the energy landscape.
	\item \textbf{Separation between magnets} $\boldsymbol{S_\mathrm{ASI}} \gtrsim \SI{20}{\nano\metre}$. \newline
	This constraint is of a practical nature, as the minimal size of lateral geometrical features is limited by the accuracy of the lithographic process. \par
	The separation only affects the strength of the magnetostatic interaction between neighbouring magnets, as approximated in \hotspice by \cref{eq:2:E_MS} and the second-order correction of \cref{eq:2:E_MS_order2}.
	In ASI, a strong MS coupling --- hence low $S_\mathrm{ASI}$ --- is often preferred.
	However, for thermally active ASI we will soon see that excessively strong MS coupling is detrimental, resulting in a ``sweet spot'' of MS coupling energies.
	\item \textbf{Number of layers} $\boldsymbol{n_{\ce{Co}}} \lesssim 8$. \newline
	This limit is of a practical nature as well --- property drift during deposition of successive layers makes it harder to achieve small $S_\mathrm{ASI}$ for increasingly tall stacks. \par
	Note that the number of layers does not affect the balance between PMA and shape anisotropy, as the number of interfaces per amount of magnetic material remains constant.
	However, since MS interactions grow $\propto n_{\ce{Co}}^2$, the inter-layer coupling --- and with it the net OOP anisotropy --- becomes increasingly significant for a large number of layers.
	\item \textbf{Diameter} $\boldsymbol{D_\mathrm{NM}} \lesssim \SI{200}{\nano\metre}$. \newline
	Magnets wider than this limit were observed to take on multi-domain states~\cite{KUR-24}, as their demagnetisation energy becomes dominant due to its rapid $\propto D_\mathrm{NM}^4$ growth.
	Even though in-plane magnets without uniform magnetisation have been used for computation~\cite{gartside2022reconfigurable}, we did not intend to use this in the OOP systems.
	Furthermore, the \hotspice simulator requires a single-domain magnetisation state for its Ising-like model to be applicable. \par
	The diameter can be used to control the relative strength between PMA and shape anisotropy and greatly affects the inter-magnet MS coupling.
	\item \textbf{Spacer layer thickness} $\boldsymbol{t_{\ce{Pt}}} \gtrsim \SI{0.4}{\nano\metre}$. \newline
	Thinner layers may become discontinuous, which is detrimental to the PMA.
	On the other hand, a thin spacer layer maximises the net OOP anisotropy due to the $\approx t_{\ce{Pt}}^{-3}$ inter-layer MS coupling dependence.
	Note that $t_{\ce{Pt}}$ must still be chosen appropriately to promote FM coupling between layers while maintaining the interfacial anisotropy. % Accounting for various effects such as RKKY coupling
\end{itemize}

In the particular case of the thermally active OOP ASI fabricated within the \spinengine project, the nanomagnets consist of $n_{\ce{Co}}=7$ ferromagnetic layers of $D_\mathrm{NM}=\SI{170}{\nano\metre}$ diameter and $t_{\ce{Co}} \approx \SI{1.45}{\nano\metre}$ thickness, with $t_{\ce{Pt}} \approx \SI{0.8}{\nano\metre}$ spacing layers in between.
The edge-to-edge spacing between neighbouring magnets is $S_\mathrm{ASL}=\SI{30}{\nano\metre}$, unless specified otherwise, with a saturation magnetisation $M_\mathrm{sat}=\SI{1063}{\kilo\ampere\per\metre}$~\cite{Msat_Co}.

\subsubsection{Relaxation characteristics}
\label{sec:3:relaxation}
To better inform our choices for RC in thermally active ASI, it is essential to familiarize ourselves with its spontaneous relaxation process.
In this section, we numerically investigate this process with \hotspice by initializing an OOP ASI in a uniform ($\uparrow$) state and subsequently observing its decay towards the checkerboard ground state over a certain timespan.
In the Ising-like model used by \hotspice, two essential system parameters remain, denoted as
\begin{itemize}[leftmargin=4.1em]
	\item[$\boldsymbol{\EEA}$ ---] the energy barrier of a single non-interacting OOP magnet due to the net OOP anisotropy.
	This is the combined contribution of PMA, shape anisotropy and the inter-layer magnetostatic interaction.
	\item[$\boldsymbol{\EMC}$ ---] the magnetostatic coupling energy between nearest neighbours.
	This is responsible for the relaxation to the ground state.
\end{itemize}
Due to the N\'eel-Arrhenius switching law~\eqref{eq:2:Néel}, it is more appropriate to write dimensionless energies ($\EEA/\kBT$ and $\EMC/\kBT$), as it is only this ratio --- not the absolute value of the energies or temperature --- that controls the switching times. \par
In real ASI, manufacturing inevitably introduces imperfections.
Slight geometrical variations can result in a significant variance of the net OOP anisotropy between magnets.
In simulations, this will be modelled by sampling $\EEA$ from a normal distribution: unless otherwise specified, a standard deviation $\sigma(\EEA) = \SI{5}{\percent}$ is used.
This provides pinning sites for domain walls, making the relaxation process more reproducible: such ``consistent variation'' within an ASI could ultimately be beneficial for RC by providing a richer output space. \par
The effective energy barrier $\EBeff$ of a magnet $i$ can then symbolically be written as
\begin{equation}
	\label{eq:3:OOP_relaxation_EBeff}
	\EBeff{}_{,i}(t) = \EEA \Big(1 + \sigma\ab(\EEA) \sampi_i \Big) + \EMC \sum_j s_i(t) D_{ij} s_j(t)
\end{equation}
with $\sampi$ a random value from a standard normal distribution and $D_{ij}$ a factor proportional to the magnetostatic interaction energy between magnets $i$ and $j$. \\\par

The relaxation process can be tracked by the normalised average magnetisation $\mavg = \abs{\llangle s \rrangle}$ and the local antiferromagnetic parameter $\qNN = (1 - \langle s_i s_{i+1} \rangle)/2$.
In the uniform state --- where all magnetic moments point `up' --- their values are $\mavg=1$ and $\qNN=0$.
During relaxation, they transition to $\mavg=0$ and $\qNN=1$, corresponding to a checkerboard ground state.
Note that, while for a random state $\qNN \approx 0.5$, a value of $\qNN=0.5$ does not necessarily imply a random state. \par

\cref{fig:3:OOP_relaxation} shows relaxation profiles for a few combinations of $\EEA$ and $\EMC$, providing some insight into the relaxation process.
Each panel shows the mean, standard deviation and 1st/99th percentile of $\mavg$ and $\qNN$ over 200 relaxations --- this variance is both due to the randomness of N\'eel-Arrhenius switching and because each relaxation used unique randomly sampled $\EEA$ for individual magnets, in accordance with the $\sigma(\EEA) = \SI{5}{\percent}$ variation.
Each relaxation consisted of at most $40N$ switches using the first-switch method, as this is sufficient for the system to achieve equilibrium.
The size of the ASI does not significantly affect the relaxation process, apart from reducing the spread on these traces.

\xfig{3_OOP/OOP_relaxation.pdf}{
	\label{fig:3:OOP_relaxation}
	Dependence of the ordering timescale on the energy balance between net OOP anisotropy $\EEA$ and nearest-neighbour magnetostatic coupling $\EMC$ in an $11 \times 11$ OOP square-lattice ASI. % This system size was chosen because it matched the size of our fabricated systems discussed later.
	The ASI is initialised in a uniform magnetic state with all spins pointing `up' and is released at $t = 0$, after which the lattice relaxes to lower energies.
	Blue traces show the absolute value of the normalised average magnetisation $\mavg$ while orange traces represent the local antiferromagnetic parameter $\qNN$.
	Each trace shows the mean (central line), standard deviation (central shaded area), and 1st/99th percentile (outer shaded area) of 200 repeated relaxations with randomly sampled $\EEA$ profiles.
	At most $40N$ switches are simulated for each relaxation.
	The green shaded region highlights the second stage of ordering.
	The simulations were stopped after \SI{e4}{\second} and the resulting final magnetic states are shown in \textbf{the insets}, where white (black) indicates up (down) spins.
}

\paragraph{Logarithmic relaxation}
Most strikingly, by using a logarithmic scale for the time axis, both $\mavg$ and (to a lesser extent) $\qNN$ trace a very straight line during relaxation.
Because of this \idx{logarithmic relaxation}, the first switches occur on a vastly different timescale than the last --- each subsequent switch occurs on average $x$ times later than the previous switch ($t_{i+1} = x t_i \Rightarrow t_n = t_0 x^n$). \par
The cause of this is the magnetostatic interaction\footnote{
	Without the \xref{magnetostatic interaction}, the system would simply be a non-interacting ensemble of nanomagnets, which is known to relax to the ground state exponentially.
	This is the opposite of logarithmic relaxation; quantities like $\mavg$ would only show up as straight lines when using a linear time axis and logarithmic y-axis.
	% Exponential decay would be much easier to work with for temporal input, but having no MS coupling defeats the point of using an ASI for some nonlinearity.
}.
Initially, all magnets are in a high-energy state as they all point in the same direction.
After a certain amount of time, a first magnet will switch, thereby slightly stabilising all other magnets in the ASI to some degree according to the second term of~\cref{eq:3:OOP_relaxation_EBeff}.
Whichever magnet randomly switches next will therefore experience a slightly higher effective energy barrier $\EBeff$ than before.
On average, this means that the effective barrier $\max(\EBeff)$ of the most unstable magnets in the system --- those most eager to switch next --- will increase linearly with the amount of switched magnets.
The N\'eel-Arrhenius equation~\eqref{eq:2:Néel} translates this to an exponentially increasing switching time. \\\par

\cref{fig:3:OOP_relaxation} also shows that the slope generally becomes less steep for higher coupling $\EMC$.
Following the previous discussion, we now understand that high coupling results in a greater decrease of $\EBeff$ per switch, and consequently a larger exponential increase in the switching time.
Meanwhile, the net OOP anisotropy $\EEA$ only affects the timescale of the relaxation --- it does not affect the slope, but simply shifts the relaxation curves along the (logarithmic) time axis.
Therefore, to achieve an ASI that exhibits a particular relaxation timescale, accurate control of $\EEA$ on a scale of $\sim 10\kBT$ is required. \par
The effects of $\EMC$ and $\EEA$ can be quantified by fitting a logarithmic curve to the mean relaxation curves of the nine panels in ~\cref{fig:3:OOP_relaxation}, yielding the general trends
\begin{equation}
	\label{eq:3:relaxation_fit_ab}
	\mavg = a \log_{10}(t \nu_0) + b \quad \mathrm{with} \quad \begin{cases}
		a \approx -\frac{\kBT}{2\EMC} \, \mathrm{,} \\
		b \approx \frac{\kBT}{6\EMC}(\frac{\EEA}{\kBT} - 1.5) - 0.05 \, \mathrm{.}
	\end{cases}
\end{equation} % TODO: more rigorous fit?

% Note that the change in slope $a$ and the offset $b$ compensate to some extent such that the final switches ($\mavg \rightarrow 0$) occur around the same time nearly independently of $\EMC$.
% The intercept of the fit for $\mavg=0$ is indeed far less sensitive to $\EMC$ than to $\EEA$:
% \begin{equation}
%  	t_{\mavg=0} = \frac{10^{-b/a}}{\nu_0} = \frac{10^{\frac{\EEA - 0.3 \EMC}{3\kBT} - 0.5}}{\nu_0} \mathrm{.}
% \end{equation}

\paragraph{Two-stage relaxation}
While $\mavg$ follows a near-perfect logarithmic relaxation, $\mavg \approx 0$ is reached earlier than the perfect checkerboard ordering $\qNN = 1$ --- most strikingly in panels 6 and 9 of~\cref{fig:3:OOP_relaxation}.
This is accompanied by a decreased ordering rate $\frac{\partial \qNN}{\partial \log(t)}$, marking the transition between stage $\circled{1}$ and $\circled{2}$ of relaxation.
The latter is indicated by the green shaded area in the figure. \par
This can be understood by considering the microstate of the ASI: to this end, the final state of each panel in~\cref{fig:3:OOP_relaxation} is shown as an inset.
Panels 1 and 2 show the system during the logarithmic relaxation of both $\mavg$ and $\qNN$ in stage $\circled{1}$.
Here, many magnets still have three or four NN with the same magnetisation state, causing random magnets to switch and small checkerboard domains to nucleate throughout the system.
Because these domains are twofold degenerate, the perfect checkerboard ordering is not reached directly.
Instead, expanding domains are likely to be separated by domain walls, as seen in panels 3 and 4 which have just barely reached stage $\circled{2}$.
From this point on, a further increase in $\qNN$ can only\footnote{
	If the magnetostatic coupling $\EMC$ is insufficient, a perfect checkerboard ordering will not be achieved and $\qNN$ may not increase much further --- in panel 7 of~\cref{fig:3:OOP_relaxation}, $\qNN$ settles around $\approx 0.8$.
} be achieved by switching magnets that are part of a domain wall.
However, even though domain wall magnets are the least stable magnets in the system, they still typically have 3 stabilising NN, resulting in a significant effective energy barrier that depends on both $\EEA$ and $\EMC$.
Furthermore, when a domain wall magnet finally switches, it simply moves the domain wall by one lattice unit without significantly altering the macrostate of the system.
This is in stark contrast with stage $\circled{1}$, where the macrostate was constantly changing.
Only when the domain wall randomly encounters a domain wall of opposite polarity or an edge of the system, can the total length of domain walls in the system decrease --- as occurs between panels 5 and 8.
This process is more akin to a random walk, and therefore has a polynomial time dependency rather than a logarithmic one.
For this reason, stage $\circled{2}$ shows up in~\cref{fig:3:OOP_relaxation} as a sudden increase of $\qNN$ after a plateau, as most clearly seen in panel 9. \par
Once the system has finally approached the checkerboard state with $\mavg \rightarrow 0$ and $\qNN \rightarrow 1$, random thermal fluctuations will still cause small excursions to $\mavg > 0$ and $\qNN < 1$.
These are generally more pronounced for low coupling $\EMC$, as can clearly be seen in the $\EEA = 20\kBT$ row of~\cref{fig:3:OOP_relaxation}.
This is because $\EMC$ is the only factor responsible for order in the system: low coupling makes magnets less inclined to reside in the ground state, giving the opportunity for multiple magnets to be in an unstable state at once and drive $\mavg$ further from 0.

%\paragraph{Miscellaneous remarks}
%During the relaxation process, one may be inclined to draw the conclusion from~\cref{fig:3:OOP_relaxation} that larger $\EMC$ yields more random variation.
%This is mostly an illusion due to the steep slope of the relaxation for low $\EMC$: while the lower slope for $\mavg$ increases the random variation along the temporal axis, the vertical spread actually decreases for larger $\EMC$. \par

\paragraph{Phase diagram}
A stronger coupling $\EMC$ causes faster decay of $\mavg \rightarrow 0$: in panels 4--6 of~\cref{fig:3:OOP_relaxation}, the time to reach $\mavg \sim 0$ decreases from $\approx \SIrange{e3}{10}{\second}$.
However, strong coupling does not necessarily result in faster ordering $\qNN \rightarrow 1$ --- as panels 8 and 9 clearly illustrate --- because domain wall magnets typically have 3 oppositely magnetised nearest neighbours, resulting in a net stabilising effect proportional to $\EMC$ (and $\EEA$).
Therefore, for a system to reach the checkerboard state on a given timescale, a sufficiently low $\EMC$ and $\EEA$ are required. \par
This is most clearly seen in phase diagrams like~\cref{fig:3:OOP_relaxation_continuous}, which summarise the values of $\mavg$ and $\qNN$ at $t = \SI{e4}{\second}$ (i.e., the maximum time in \cref{fig:3:OOP_relaxation}) over a wide range of $\EEA$ and $\EMC$.
This encompasses the nine $(\EMC,\EEA)$ combinations of the various panels in~\cref{fig:3:OOP_relaxation}, as indicated by grey numbered crosses in~\cref{fig:3:OOP_relaxation_continuous}. \par
Five distinct regions can be distinguished.
In region $\mathrm{I}$ (panel 1), weak coupling and high anisotropy result in a non-volatile system --- as used in the previous section --- where no spontaneous switches occur and $\mavg = 1$ and $\qNN = 0$ remain unchanged.
Increasing the $\EMC/\EEA$ ratio shifts the system though the transient region $\mathrm{II}$ (panel 2), characterised by decreasing $\mavg$ and increased ordering.
Further increase of $\EMC/\EEA$ shifts the system to region $\mathrm{III}$ (panel 6) with $\mavg \approx 0$ but $\qNN \neq 1$, characteristic of stage $\circled{2}$.
Any further increase of $\EMC/\EEA$ does not affect the average magnetisation or ordering at $t = \SI{e4}{\second}$ --- unless both $\EEA$ and $\EMC$ are sufficiently small, as in regions $\mathrm{IV}$ and $\mathrm{V}$.

\xfig{3_OOP/OOP_relaxation_continuous.pdf}{
	\label{fig:3:OOP_relaxation_continuous}
	Phase diagrams of the normalised average magnetisation $\mavg$ and the local antiferromagnetic parameter $\qNN$ as a function of net out-of-plane anisotropy $\EEA$ and nearest-neighbour magnetostatic coupling $\EMC$ at $t=\SI{e4}{\second}$.
	Five regions can be distinguished: $\mathrm{I}$ -- frozen state, $\mathrm{II}$ -- transient states, $\mathrm{III}$ -- state with domains and domain boundaries, $\mathrm{IV}$ -- checkerboard ordering, $\mathrm{V}$ -- superparamagnetic state.
	The white and black dotted line respectively correspond to $\mavg = 0.05$ and $\qNN = 0.95$.
	The nine panels of~\cref{fig:3:OOP_relaxation} are indicated by grey crosses.
	The three symbols correspond to the experimental lattices shown in~\cref{fig:3:OOP_relaxation_fit} with nanomagnet diameter $D_\mathrm{NM} = \SI{170}{\nano\metre}$, \ce{Co} thickness $t_{\ce{Co}} = \SI{1.45}{\nano\metre}$ and nanomagnet separations $S_\mathrm{ASI} = \SI{20}{\nano\metre}$ (circle), $S_\mathrm{ASI} = \SI{25}{\nano\metre}$ (triangle), $S_\mathrm{ASI} = \SI{30}{\nano\metre}$ (cross). % TODO: create and reference that figure
}

Region $\mathrm{IV}$ (panel 9) is the only region where the system spontaneously achieves the checkerboard ground state.
In~\cref{fig:3:OOP_relaxation_continuous}, it is delineated by the dotted black line that envelops a region with $\qNN \geq 0.95$.
Reducing $\EMC$ too much puts the system in region $\mathrm{V}$ (panel 7), where the system is superparamagnetic at this timescale.
For $t = \SI{e4}{\second}$, region $\mathrm{V}$ extends upwards to $\EEA \lesssim 32 \kBT$, since the N\'eel-Arrhenius switching law states that nanomagnets with such anisotropy have a mean switching time $\lesssim \exp(32)/\nu_0 \approx \SI{e4}{\second}$.
The main region of interest, region $\mathrm{IV}$ with $\qNN \approx 1$, is therefore constrained from the left by the superparamagnetic regime, from the top by the transient/frozen regime, and from the right by states in which domain boundaries are too stable at the given timescale. \par
The phase diagrams in~\cref{fig:3:OOP_relaxation_continuous} are snapshots at $t = \SI{e4}{\second}$ and, with increasing time, region $\mathrm{IV}$ will expand and the boundaries between the different regions will shift upwards.
Therefore, if the observation is long enough and $\EMC$ is non-negligible, a system located in the frozen region at $t = \SI{e4}{\second}$ may eventually find itself in the more ordered Regions $\mathrm{II}$, $\mathrm{III}$ or $\mathrm{IV}$.
However, due the exponential character of the N\'eel-Arrhenius law, this shift will be minimal and region $\mathrm{IV}$ will not become much bigger than in~\cref{fig:3:OOP_relaxation_continuous} for any realistic timescale~\cite{KUR-24}.

\paragraph{Exchange-coupled system} % TODO: expand, could be omitted if not sufficiently worked out

\paragraph{Summary}
The state of an OOP ASI is defined by the net OOP anisotropy $\EEA$, magnetostatic coupling $\EMC$ and elapsed time $t$.
Increasing $\EMC$ causes the relaxation to occur over a larger amount of timescales.
Meanwhile, $\EEA$ controls the timescale of the relaxation but does not affect the amount of timescales over which the relaxation occurs.
For reservoir computing purposes, a sufficiently low $\EEA$ and $\EMC$ are required --- but not too low, to maintain OOP anisotropy and a ground-state checkerboard ordering, respectively.
Therefore, practical applications exploiting systems with specific dynamics will require careful adjustment of the lattice and nanomagnet dimensions~\cite{KUR-24}.

\subsubsection{Fitting MFM images}\label{sec:3:OOP:MFM}
Based only on the magnet geometry --- determined by the various geometrical parameters discussed earlier in~\cref{sec:3:OOP_geometry} --- it is hard to predict the net OOP anisotropy $\EEA$ and magnetostatic coupling $\EMC$ of a manufactured ASI.
For typical nanomagnet dimensions (7 layers of $t_{\ce{Co}} \approx \SI{1.45}{\nano\metre}$ with $D_\mathrm{NM} \approx \SI{170}{\nano\metre}$), some of the aforementioned energy contributions are on the order of several $1000 \kBT$, as we discuss in more detail in~\ccite{KUR-24} (Supp. 4). % --- especially PMA and shape anisotropy, which have opposite sign and therefore mostly cancel out.
Due to manufacturing defects, the discrepancy between theoretical estimates of these energy contributions and their real magnitudes can easily exceed $100 \kBT$. % Theoretically, $\EEA$ of an isolated magnet (as manufactured by Alex and observed to be part of a thermally relaxing OOP ASI) is $\approx\SI{-3.6e-18}{\joule} = \SI{-22.5}{\electronvolt} = -871\kBT$ which is far too large to get thermal switching. Furthermore, it is negative, implying an IP magnetisation. This discrepancy between Mathematica and reality could be due to non-coherent reversal, lowering the real energy barrier.
However, we just saw that $\EEA$ and $\EMC$ must not exceed several tens of $\kBT$ if a system is to spontaneously reach the checkerboard ground state.
It is therefore inappropriate to calculate $\EEA$ and $\EMC$ theoretically --- they must be empirically determined from manufactured ASI. \\\par
For this reason, \textit{A. Kurenkov} manufactured several $11 \times 11$ OOP square-lattice ASI, initialised them in the `up' state by a \SI{0.9}{\tesla} OOP field and measured their state using MFM at different time intervals since initialisation.
For each of these MFM images, the values of $\mavg$ and $\qNN$ can be determined.
The knowledge of these quantities at different moments in time can then be used to fit this data to simulated relaxations, like the ones in~\cref{fig:3:OOP_relaxation}, by varying simulation parameters like $\EEA$ and $\EMC$.

\paragraph{Bayesian optimisation}
Any fitting procedure requires an appropriate metric $\eta$ to determine the quality of the fit.
For this, we will use the random distribution of both $\mavg$ and $\qNN$ at the times of measurement $t_i$, as previously seen in~\cref{fig:3:OOP_relaxation}.
These distributions are each characterised by a mean and standard deviation, which can be determined from repeated simulations --- typically, 200 relaxations are used. % We embrace the Monte Carlo nature of \hotspice.
Therefore, each experimental data-point lies some amount of standard deviations away from the mean of this simulated distribution.
Assuming a normal distribution\footnote{
	Since the random relaxation is governed by the N\'eel-Arrhenius equation~\eqref{eq:2:Néel}, quantities like $\mavg$ which are bounded in the range $[0,1]$ most likely follow a binomial distribution, which can be well approximated by a normal distribution for a large number of magnets $N$ in the system (here: $N > 100$).
}, a probability density can be assigned to this.
As fit metric $\eta$, we will use the product of these probability densities, multiplied over each experimental $\mavg$ and $\qNN$ at each measured time $t$.
This can be interpreted as a relative probability that the MFM measurement can be explained by a simulation using certain system parameters.
Mathematically,
\begin{equation}
	\eta = \bigprod_{t_i} \bigprod_{x \in \left\{\substack{\mavg\\\qNN\\\qNN[2]\\\qNN[3]}\right.} \exp\ab[-\frac{1}{2} \ab(\frac{\langle x_\mathrm{sim}(t_i) \rangle - x_\mathrm{MFM}(t_i)}{\sigma\big(x_\mathrm{sim}(t_i)\big)})^2] \mathrm{,}
\end{equation}
where a superscript `sim' indicates the simulated value, while `MFM' indicates the measured value from MFM images taken at various times $t_i$.
For a more accurate fit, we also used the longer-range order parameters $\qNN[2]$ and $\qNN[3]$, defined similar to $\qNN = (1 - \langle s_i s_{i+1} \rangle)/2$ but instead with a correlation between next-nearest neighbours\footnote{
	2NN in OOP square ASI are diagonally adjacent magnets.
} and third-nearest neighbours\footnote{
	3NN in OOP square ASI are magnets two lattice units apart horizontally or vertically.
}. \par
With this\footnote{
	Because $\eta$ can easily vary over many orders of magnitude, the optimiser will instead use $\log_{10}(\eta)$.
} fit metric $\eta$, the optimum combination of $\EEA$ and $\EMC$ can be determined.
We will also allow the exchange coupling $J$ to be varied during fitting --- if the fit yields a significantly non-zero value of $J$, this can indicate that the system contains defects in the form of interconnected magnets. \\\par

The fit metric $\eta$ is rather expensive to compute, because calculating the distribution of relaxations takes a non-negligible amount of time.
Therefore, an optimisation method must be used that can find the optimum with a limited number of $\eta$-evaluations.
\idx{Bayesian optimisation} (BO)~\cite{bayesopt_package,BayesOpt_Mockus1975} is one such method that can find the extremum of an expensive-to-compute function in a low-dimensional space (here, we only search the $\EEA \times \EMC \times J$-space).
This is an iterative method that starts from an initial guess --- estimated values of $\EEA$ and $\EMC$ can be determined based on the experimentally measured $\mavg(t)$ by using the inverse of~\cref{eq:3:relaxation_fit_ab}.
It then explores a region in $\EEA \times \EMC \times J$-space by choosing new parameter combinations based on an acquisition function --- here, the upper confidence bound is used.
While BO may not be the most performant method, as it converges rather slowly and its computational complexity scales with the number of iterations $n$ as $\order{n^3}$, it is sufficient for our purposes.
More extensive parameter searches may benefit from using other optimisation techniques like LIPO~\cite{LIPO,LIPO_dlib}. \\\par % Good explanation in http://blog.dlib.net/2017/12/a-global-optimization-algorithm-worth.html

Due to the logarithmic relaxation, an accurate fit requires measurements separated by multiple orders of magnitude in time.
Because $\mavg \approx a \log_{10}(t \nu_0) + b$, two measurements should suffice to fit this kind of curve to experimental data.
While two measurements are insufficient to prove that the relaxation indeed follows a logarithmic behaviour as predicted by simulation, this was verified using multiple measurements over a shorter timescale as discussed in~\ccite{KUR-24} (Supp. 5.4).
For this reason, the samples were measured soon ($\approx \SI{e3}{\second}$) after initialisation, and again after 2 years $\approx \SI{7e7}{\second}$, providing two sufficiently distinct timescales for a successful fit.

\paragraph{Results} % EEA etc.
% TODO: present the results of this BO procedure, possibly in combination with the `simple' vs. `parabolic' energy barrier and the infinitesimal spin vs. the finite spin approximation in the DD energy calculation, and how these affect the best fit (and coupling this back to verifying which approximations are beneficial and which are not).

\subsection{RC in thermally active OOP square ASI}
\subsubsection{Concept} % Lessons learnt from non-volatile ASI
\paragraph{Input torques: representation in \hotspice}
Contrary to the binary input in the non-volatile OOP ASI, thermally active ASI can be addressed by a continuous range of input.
The input --- in reality applied through spin-orbit torque --- was applied in the simulations in the form of an OOP magnetic field $B_z(t)$ acting on the entire system.
The field was scaled such that its OOP component extends from $B_0$ (minimum) to $B_1$ (maximum), i.e. $B_0 \leq B_z \leq B_1$.

\paragraph{Readout method}
The inverse mean squared error (1/MSE) between o(t) and the desired result (the sawtooth) was used as a performance metric.

\subsubsection{RC: signal transformation}
\paragraph{Motivation} % Why do signal transformation? (scalar, nonlinear, memory...)
\paragraph{Baseline}
\paragraph{Parameter dependencies}
\subparagraph{Field magnitude}
\subparagraph{Frequency}
\subparagraph{Gradient}
\subparagraph{System size and readout resolution}
\subparagraph{Exchange/no exchange}
\subparagraph{Transformation}
\paragraph{Bayesian optimisation: Mackey-Glass}
\paragraph{Reflecting on results} % Is the performance decent? Compare to literature. MSE is not an ideal parameter, should use something that accounts more for the general shape.
\cleardoublepage
\chapter{Conclusion and outlook}\label{ch:Conclusion}
\glijbaantje{There is no prize to perfection, only an end to pursuit.}{Viktor, \textit{Arcane}}

Throughout this thesis, we have investigated the use of \xref{artificial spin ice} with perpendicular magnetisation --- of both the thermally active and non-volatile kind --- as a physical \link{reservoir computing}{reservoir}.
Our interest in OOP ASI for \xref{reservoir computing} was motivated by its well-established input (\link{spin--orbit torque}{SOT}) and readout (\link{anomalous Hall effect}{AHE}) methods, but its RC potential had yet to be quantified.
To support this work, we developed \hotspice, a flexible Monte Carlo simulator that captures the collective behaviour of large ASI arrays.
Altogether, we have covered three main topics.
\begin{itemize}
	\item The development of \textbf{\hotspice}, with a focus on several improvements that increase the physical accuracy and algorithmic performance of the underlying model.
	\item The demonstration that \textbf{thermally active} OOP ASI can serve as an analog-input reservoir by harnessing its \link{logarithmic relaxation}{spontaneous relaxation}, as confirmed through non-linear signal transformation and prediction.
	\item The design and numerical validation of a \link{AFM clocking protocol}{clocking protocol} for \textbf{non-volatile} OOP ASI that provides \link{stepwise domain wall motion}{controlled domain wall motion}, enabling it to be used as a binary-input reservoir.
\end{itemize}
The key conclusions are summarised below, followed by an outline of promising avenues for future research.

\newpage
\section{Summary of main results}
The \link{point dipole model}{Ising-like model} used by \textbf{\hotspice} to represent \xref{single-domain nanomagnets} in an ASI allows it to strike a practical balance between physical accuracy and computational efficiency.
Two distinct switching algorithms were implemented: the \xref{first-switch method} is best used to simulate temporal dynamics, while \link{Metropolis-Hastings sampling}{Metropolis-Hastings (MH) sampling} efficiently explores the equilibrium state space.
The latter can gain performance by simultaneously \link{multi-switching}{sampling multiple} (sufficiently distant) magnets.
Two classes of model improvements were found to enhance simulation accuracy at little to no performance impact.
Firstly, accounting for the finite size of magnets in OOP ASI is best done through a \link{finite dipole model}{second-order correction} on the MS interaction.
For IP ASI, on the other hand, a \xref{dumbbell model} is preferable and can yield qualitatively different dynamics in, e.g., kagome ASI.
Secondly, to accurately estimate the \xref{effective energy barrier} in IP ASI, accounting for \link{asymmetric energy barrier}{asymmetric switching channels} and \link{non-coherent magnetisation reversal}{non-coherent reversal} proved key to reproduce experimental results.
For OOP magnets, however, an exact solution for an idealised energy landscape proved inaccurate.
In both ASI types, the assumption of switching by \link{coherent rotation}{coherent reversal} therefore seems to be a poor approximation. \par
Using \hotspice, it was shown that \textbf{thermally active} OOP ASI exhibits \xref{logarithmic relaxation} towards its ground state, a process characterised in detail in~\cref{sec:3:relaxation}.
External stimuli can push it away from equilibrium, leading to behaviour reminiscent of a \xref{leaky integrator} over a remarkably broad frequency window ($\gtrsim 2$ decades) owing to the logarithmic relaxation.
This allows it to serve as a \link{fading memory}{fading-memory} \link{reservoir computing}{reservoir}, as verified by two benchmark tasks --- non-linear signal transformation and \link{Mackey-Glass oscillator}{chaotic time-series} prediction --- where performance comparable to other physical reservoirs was obtained ($\MSE \lesssim 0.015$).
Including spatial variations, such as a \link{property gradient}{gradient in material properties}, was found to have a positive impact as this introduces different relaxation timescales throughout the system, thereby encoding more distinct information. \par
For \textbf{non-volatile} OOP ASI, global stimuli would saturate the system, as \xref{avalanches} can propagate unimpeded --- this was not an issue for thermally active ASI, since its finite relaxation time allows it to gradually respond to input.
Therefore, we developed the \link{AFM clocking protocol}{``AFM clocking'' protocol}: by alternately stimulating two interleaved sublattices, avalanches were avoided and controlled \xref{stepwise domain wall motion} was achieved.
Defects, in the form of a variation on the OOP anisotropy, enable this scheme to produce more complex dynamics that can be harnessed for RC, as the position and size of domains non-linearly encodes recent bits of information.
In general, we can say that controlled domain wall motion in any ASI requires, at the very least, (1) an input that lifts the ground state degeneracy and (2) at least two independently addressable sublattices, to prevent avalanches. \\\par

In summary, we have established that OOP ASI can indeed function as a versatile \link{physical reservoir computing}{physical reservoir}: in its thermally active form, it responds to analog input akin to a \xref{leaky integrator}, while its non-volatile counterpart supports binary input through the \link{AFM clocking protocol}{clocking protocol}.
Meanwhile, \hotspice has proven its worth as a flexible and efficient simulator, usable for both the exploration of ASI physics and their optimisation in \link{neuromorphic computing}{unconventional computing} architectures.

\newpage
\section{Outlook and future directions}
Looking ahead, several avenues can strengthen both the \hotspice simulation toolkit and the proposed \link{physical reservoir computing}{physical reservoir} concepts.

\paragraph{Improvements to \hotspice}
\hotspice has gradually grown as new improvements to physical accuracy or performance were consecutively introduced, which has led to some features having a rather opaque API.
Now that the main functionality is well-defined and has been validated, the public API should be streamlined by removing deprecated functions and clearly exposing the most commonly used routines. % while retaining the versatility of the current implementation.
This will make the package easier to maintain and more approachable for prospective users.
While the various implemented improvements to the \link{point dipole model}{underlying model} already make \hotspice usable in a \link{fig:2:ASIs}{wide range of lattices}, several other factors listed below could be kept in mind during refactoring, to raise the efficiency of \hotspice to the next level. \\\par

Several algorithmic improvements remain possible.
First, for equilibrium studies, the usage of \link{cluster algorithms}{clustering algorithms} (e.g., \xref{Wolff}) could drastically reduce \xref{critical slowing down}, though this route was previously abandoned since it is non-trivial to correctly apply such algorithms to the broad range of ASI lattices that \hotspice aims to simulate (i.e., arbitrarily oriented magnets with full \link{magnetostatic interaction}{dipolar coupling} between them).
Second, the \xref{first-switch method} is limited by the fastest-switching magnets in the system, which can pose an issue in situations where only a small set of magnets repeatedly switch back-and-forth.
Due to its single-switch nature, its application is also restricted to small systems: a method akin to MH \xref{multi-switching} that alleviates these restrictions would be a welcome addition. \par
We chose to represent ASI on a \xref{rectilinear grid} in \hotspice, trading some geometrical freedom for enhanced performance. %, though many popular lattices can nonetheless be recreated.
Most notably, the grid enables the efficient calculation of \textit{all} \xref{magnetostatic interactions} by means of convolutions, whereas other codes are often forced to truncate the interaction distance.
We found that such \link{truncated kernel}{truncation} incurs a non-negligible error on the total magnetostatic energy --- exceeding $\SIrange{5}{10}{\percent}$ even when accounting for the $\approx 1000$ nearest magnets --- though it can be argued that the \xref{macrospin approximation} itself constitutes a more severe approximation than this, possibly rendering such truncation justifiable.
Besides this, the grid also enables efficient \link{multi-switching}{multi-sampling} of sufficiently distant magnets for the MH algorithm, which would be much harder on an unstructured lattice. \par
For the relatively small lattices considered in~\cref{ch:Applications}, it was found that CPUs outperform GPUs at the array operations underlying this grid. % except in~\cref{fig:3:Clocking_massive_seeded}
Nonetheless, the effort of porting \hotspice to GPU was not wasted, as it has expanded the scope of \hotspice by allowing it to handle large ASIs.
Looking forward, one remaining performance hurdle is the processing of many small simulations --- like parameter sweeps or determining ensemble averages --- which currently requires a batch of separate simulations to be run.
It would be more efficient to harness the parallel processing capabilities of the GPU by instead ``stacking'' multiple simulations as 3D arrays.
While this would require significant refactoring of \hotspice, it could be a worthwhile endeavour to obtain significantly higher throughput in large ensembles of small ASIs.

\paragraph{Reservoir computing in perpendicular-anisotropy artificial spin ice}
We have successfully used both thermally active and non-volatile out-of-plane ASI as a \link{reservoir computing}{reservoir}, albeit with very different input schemes.
While the presented approaches complement each other in some ways, they both have their limits that can affect the viability of their use as \link{physical reservoir computing}{physical reservoirs}. \\\par

Generally, OOP ASI presents less design freedom than IP ASI, as the \xref{easy axis} of each \link{single-domain nanomagnet}{magnet} is always vertical.
We therefore only considered the square lattice, as an example of a lattice without nearest-neighbour \xref{frustration}.
It can be interesting to consider whether frustrated OOP lattices (e.g., triangle, Cairo\dots) might exhibit more chaotic dynamics --- though the reservoir protocols used here would likely yield very similar results regardless of the precise lattice used because all OOP lattices have a more or less \link{antiferromagnetic}{AFM} ground state.
As such, it may be more worthwhile to explore other approaches that introduce spatial variations, akin to the \xref{property gradient}, like uneven current distributions or local input. \\\par

While a \link{AFM clocking protocol}{clocking scheme} in non-volatile ASI enables \link{stepwise domain wall motion}{controlled dynamics}, it is limited to discrete input and requires intricate manufacturing of conduction pathways.
However, analog data is more prevalent in nature and allows a higher information density, so it is likely desirable to use the full potential of analog input whenever possible.
A few approaches were proposed to combine the clocking scheme with analog input, but were ultimately dismissed for not yielding the rich dynamics we require.
Thermally active ASI, on the other hand, can readily act as a \xref{leaky integrator} of analog input. \par
Manufacturing both non-volatile and thermally active ASI reservoirs is not so straightforward, but both for different reasons.
Non-volatile ASI requires more intricate layouts of current lines for input, which must avoid electrical conflicts with the readout method, whereas for thermally active ASI a full conductive underlayer is sufficient to use it as a reservoir.
Conversely, thermally active nanomagnets pose greater fabrication challenges, as precise control over the \xref{effective energy barrier} is required --- on the order of tens of $\kBT$. \par
In this regard, it is important to note that we have encountered indications of \link{non-coherent magnetisation reversal}{non-coherent reversal} in both IP (\cref{sec:2:Applications_reversal_Pinwheel}) and OOP (\cref{sec:3:MFM}) ASI, in the form of a lower energy barrier than would be expected for \link{coherent rotation}{coherent reversal}.
For the OOP magnets, it is therefore unclear whether downscaling would aid to achieve the required level of control over the effective energy barrier.
A better understanding of the reversal process may be achieved through \link{micromagnetic theory}{micromagnetic simulations}.
Overall, the experimental aspect of these approaches requires further consideration. \\\par

Regarding the broader perspective of \xref{physical reservoir computing} (PRC), many physical systems fulfil the theoretical requirements of a reservoir and can therefore act as universal function approximators, but their practical use remains limited.
%Various factors contribute to the difficulty of applying PRC in real-world scenarios.
Unlike software reservoirs, physical reservoirs must obey physical principles, which restricts their tunability and complicates scaling.
This was also apparent in OOP ASI, where enlarging the system did not yield proportional gains in performance.
We also encountered the trade-off between memory and non-linearity: reconfigurable reservoirs provide an interesting avenue in this regard, as they can dynamically balance this trade-off depending on the task at hand.
For instance, the non-volatile OOP ASI could do this by changing the input magnitude. \par
Nonetheless, physical reservoirs often suffer from limited memory capacity, as past excitations are gradually overwritten by new input.
Pre-processing of the raw input data into a more digestible form is therefore often required to excite desirable dynamics in the physical substrate.
A typical example of this is speech recognition, where a frequency spectrum of the spoken word provides a more digestible input than the time-domain audio signal, which ultimately benefits the reservoir performance as it eases the memory requirements.
It would be beneficial if such pre-processing could also be realised physically rather than in software.
It is also possible to enhance performance by embedding the reservoir into higher-level paradigms such as the single dynamical node~\cite{appeltant2011information} or rotating neurons reservoir~\cite{RotatingNeuronsRC}, though at the cost of reintroducing abstractions into the architecture. \par
I believe the recent review by Yan~\etal{}~\cite{ChallengesFutureRC} paints a realistic picture of the future role of PRC.
Physical reservoirs will not replace conventional computers, nor should they.
Instead, both systems can complement one another, with PRC being employed at the edge to provide early low-power processing of incoming data inside a larger architecture.
PRC is part of the more general idea to use the inherent processing power of physical systems when possible --- a concept which certainly has its merit but can be hard to reconcile with current computing architectures. \\\par

In conclusion, while OOP ASI is alluring for its ease of electrical interfacing, the aforementioned factors complicate its usage as a \link{physical reservoir computing}{physical reservoir}.
Nonetheless, we succeeded in using it as such, with a performance on par with other types of reservoirs described in literature.
Following this numerical exploration of the RC potential of OOP ASI, a lithographic experimental realisation of the \link{AFM clocking protocol}{clocking protocol} would be desirable, to achieve a better understanding of its performance, feasibility and efficiency in a real environment.


% Furthermore, for the functionality they enable, they take up a non-negligible amount of space on-chip as compared to traditional computer architectures, as the magnets have a size on the order of $\SI{100}{\nano\metre}$ (it is indeed possible to make the magnets smaller, as a high magnetostatic coupling is not required in thermally active ASI, but controlling the OOP anisotropy might be hard), a space in which typically one transistor can be implemented including peripherals like current lines, though one could argue that (thermally active) ASI can react to analog input signals whereas representing floating-point values would require a large number of transistors.

%It was found that decent values of the computing capacity $C$ do not necessarily translate to significantly non-zero memory capacity or non-linearity.
\cleardoublepage

\addtocontents{toc}{\protect\setcounter{tocdepth}{0}}
\appendix

\backmatter
%\bibliographystyle{bib/hunsrtnat}
\bibliographystyle{bib/unsrtnat2}
\bibliography{bib/bibliography.bib} % Formatting bibliography.bib: use https://flamingtempura.github.io/bibtex-tidy/index.html?opt=%7B%22modify%22%3Atrue%2C%22curly%22%3Atrue%2C%22numeric%22%3Atrue%2C%22months%22%3Afalse%2C%22space%22%3A2%2C%22tab%22%3Atrue%2C%22align%22%3A14%2C%22blankLines%22%3Atrue%2C%22sort%22%3A%5B%22title%22%5D%2C%22duplicates%22%3A%5B%22key%22%2C%22doi%22%5D%2C%22stripEnclosingBraces%22%3Afalse%2C%22dropAllCaps%22%3Afalse%2C%22escape%22%3Afalse%2C%22sortFields%22%3A%5B%22title%22%2C%22shorttitle%22%2C%22author%22%2C%22year%22%2C%22month%22%2C%22day%22%2C%22journal%22%2C%22booktitle%22%2C%22location%22%2C%22on%22%2C%22publisher%22%2C%22address%22%2C%22series%22%2C%22volume%22%2C%22number%22%2C%22pages%22%2C%22doi%22%2C%22isbn%22%2C%22issn%22%2C%22url%22%2C%22urldate%22%2C%22copyright%22%2C%22category%22%2C%22note%22%2C%22metadata%22%5D%2C%22stripComments%22%3Afalse%2C%22trailingCommas%22%3Afalse%2C%22encodeUrls%22%3Afalse%2C%22tidyComments%22%3Atrue%2C%22removeEmptyFields%22%3Afalse%2C%22removeDuplicateFields%22%3Afalse%2C%22lowercase%22%3Atrue%2C%22backup%22%3Atrue%7D
\cleardoublepage

\chapter{List of constants and symbols}  % TODO END: put all symbols in this list
\section*{List of constants}
\begin{longtable}[l]{p{60pt} p{140pt} p{200pt}}
	$e$ & Electron charge & \SI{1.6022e-19}{\coulomb} \\
	% $\gamma$&gyromagnetic ratio&$1.76086\times 10^{11}$\,rad/Ts\\
	% $\gamma_0$&$\mu_0\gamma$&$2.21\times\,10^5$\,m/As\\
	$k_\mathrm{B}$ & Boltzmann constant & \SI{1.3087e-23}{\joule\per\kelvin} \\
	$\mu_0$ & Vacuum permeability & $4 \pi \times \SI{e-7}{\tesla\metre\per\ampere}$ \\
	% $\mu_\mathrm{B}$&Bohr magneton&$9.274\times 10^{-24}$\,Am$^2$\\
\end{longtable}

\section*{List of symbols}
\subsection*{Mathematical operators}
\begin{longtable}[l]{p{60pt} p{350pt}}
	$\langle \cdot \rangle$ & Spatial average over the ASI \\
	$\langle \cdot \cdot \rangle$ & Correlation over the ASI \\
	$\odot$ & Hadamard product (element-wise matrix multiplication) \\
	$*$ & Convolution \\
	% $\overline{a}$& time average\\
	%$\av{a b}$& correlation\\
	%$\sav{a}$&spatial average over computational domain\\
	%$\dot{a}$&time derivative\\
	$K()$ & Complete elliptic integral of the first kind \\
	$\sigma( \cdot )$ & Standard deviation for variable sampled from Gaussian distribution \\
	$\norm{ \cdot }$ & Vector norm\\
	&\\
\end{longtable}

\subsection*{Roman symbols}
\begin{longtable}[l]{p{60pt} p{350pt}}
	$a$ & Semi-major axis ($l/2$) \\
	&\\

	$b$ & Semi-minor axis ($w/2$) \\
	$\vc{B}_\mathrm{ext}$ & External magnetic field \\
	&\\

	$d$ & Effective distance between north and south poles of a magnet in the dumbbell approximation \\
	&\\

	$E$ & Total interaction energy \\
	$\EB$ & Energy barrier due to shape anisotropy \\ % TODO: what to do with E_B^c? Make a new global rule that (something)^c is critical? Should change T_c then
	$\EBeff$ & Effective energy barrier \\
	$E_\mathrm{Exch}$ & Exchange energy\\
	$E_\mathrm{MS}$ & Magnetostatic interaction energy \\
	$E_\mathrm{Z}$ & Zeeman energy\\
	$E_\perp$ & Energy of a magnet when pointing \ang{90} counter-clockwise from its current state \\
	$\Delta E$ & Switching energy \\
	&\\

	$\mathcal{I}$ & Geometrical factor used for second-order correction on magnetostatic interaction \\
	&\\

	$J$ & Exchange coupling constant \\
	&\\

	$\vc{K}_i$ & Demagnetisation kernel for a magnet $i$ \\
	$\vc{\mathcal{K}}^{(q)}$ & Demagnetisation kernel for a magnet at site $q$ in the unit cell \\
	$K_u$ & Uniaxial shape anisotropy \\
	&\\

	$l$ & Length of elliptical magnet ($2a$) \\
	$L_x$ & Number of cells in the simulation along the x-direction \\
	$L_y$ & Number of cells in the simulation along the y-direction \\
	&\\

	$M$ & Average magnetization of ASI \\
	$M_\parallel$ & Average magnetisation of ASI along the direction of the external field $\vc{B}_\mathrm{ext}$ \\
	$\vc{M}$ & Magnetisation \\
	$M_\mathrm{sat}$ & Saturation magnetisation \\
	&\\

	$N$ & Number of magnets in the ASI \\ % TODO: add N_x and N_y?
	$\mathcal{N}_i$ & Nearest neighbours of magnet $i$ \\
	&\\

	$q$ & Magnetic charge \\
	$Q$ & Maximum allowed change in switching probability for two simultaneously sampled magnets \\
	&\\

	$\vc{r}_i$ & Position of magnet $i$ \\
	$r_{ij}$ & Distance between magnets $i$ and $j$ \\
	$\vc{r}_{ij}$ & Distance vector between magnets $i$ and $j$ \\
	% TODO: R_x and R_y? Are grid-select algorithm symbols so might not be very useful to mention.
	$\vc{r_N}_i$ & Position of north monopole of magnet $i$ in the dumbbell model \\
	$\vc{r_S}_i$ & Position of south monopole of magnet $i$ in the dumbbell model \\
	&\\

	$s_i$ & State of the macro-spin along the easy axis ($\pm 1$) \\ % TODO: better wording
	$\vc{s}$ & Column vector containing all $s_i$ \\
	$\langle S_i S_{i+1} \rangle$ & Nearest-neighbour correlation \\
	&\\

	$t$ & Time \\
	$\Delta t$ & Random switching time interval of a single magnet \\
	$T$ & Temperature \\
	$T_c$ & Critical temperature (Ising model) \\
	&\\

	$\vc{u}_i$ & Uniaxial anisotropy axis of magnet $i$ \\
	&\\

	$V$ & Volume of a magnet \\
	&\\

	$w$ & Width of elliptical magnet ($2b$) \\
	&\\
\end{longtable}

\subsection*{Greek symbols}
\begin{longtable}[l]{p{60pt} p{350pt}}
	% $\alpha$&Gilbert damping parameter\\
	% &\\

	% $\beta$&$=\frac{KV}{k_\mathrm{B} T}$, anisotropy energy with respect to thermal energy\\
	% &\\

	% $\gamma$& gyromagnetic ratio\\
	% $\gamma_0$& $\mu_0\gamma$\\
	% $\Gamma_{XY}$&coviariance between stochastic variables $X$ and $Y$\\
	% &\\

	$\delta$ & Relative magnetostatic/exchange coupling \\
	&\\

	% $\eta$& viscosity\\
	% &\\

	$\theta$ & In-plane magnetization angle measured clockwise from the easy axis $\vc{u}$ \\
	&\\

	% $\kappa$& windowing weighting factor\\
	% &\\

	% $\lambda$&index parameter to denote a single event or realization of a stochastic variable\\
	% &\\


	$\mu_i$ & Size of magnetic moment of magnet $i$ \\
	$\vc{\mu}_i$ & Total magnetic moment vector of magnet $i$ \\
	&\\

	
	$\nu$ & Switching rate \\
	$\nu_0$ & Attempt frequency \\
	&\\

	$\rho$ & Perpendicular magnetization factor \\
	&\\

	% $\sigma$&standard deviation in the lognormal distribution\\
	% $\sigma_X$&standard deviation of the stochastic variable $X$\\
	% $\sigma_X^2$&second central moment of stochastic variable $X$, variance\\
	% &\\

	% $\tau_0$& $\frac{1}{2\nu_0}$\\
	% $\tau_\mathrm{B}$& Brownian fluctuation time constant\\
	% $\tau_\mathrm{eff}$& effective fluctuation time constant\\
	% $\tau_\mathrm{ex}$& timescale of the experiment\\
	% $\tau_\mathrm{N}$& N\'eel fluctuation time constant\\
	% &\\

	$\phi$ & Angle of the effective field w.r.t. easy axis $\vc{u}$ \\
	&\\

	$\chi$ & Uniformly distributed random variable in $(0, 1]$ \\
	&\\
\end{longtable}
\cleardoublepage
\chapter{List of abbreviations}
\markboth{List of abbreviations}{List of abbreviations} % For headers in fncychap
{
	\addtolength{\skip\footins}{1pc}
    \begin{longtable}[l]{ll}
        AFM   & \link{antiferromagnetic}{Antiferromagnetic} \\
        AHE   & \link{anomalous Hall effect}{Anomalous Hall effect} \\
        AMR   & \link{anisotropic magnetoresistance}{Anisotropic magnetoresistance} \\
        ANN   & \link{artificial neural network}{Artificial neural network} \\
        ASI   & \link{artificial spin ice}{Artificial spin ice} \\
        BC    & \link{periodic boundary conditions}{Boundary conditions} \\
        BO    & \xref{Bayesian optimisation} \\
        BPTT  & \link{backpropagation through time}{Backpropagation through time} \\
        CPU   & \link{central processing unit}{Central processing unit} \\
        CSD   & \link{critical slowing down}{Critical slowing down} \\
        EA    & \link{net out-of-plane anisotropy}{Effective anisotropy} \\
        ESN   & \link{echo state network}{Echo state network} \\
        % ETHZ  & Eidgen\"ossische Technische Hochschule Z\"urich \\
        FD    & \link{finite-difference}{Finite-difference} \\
        FMR   & \link{ferromagnetic resonance}{Ferromagnetic resonance} \\
        FNN   & \link{feed-forward neural network}{Feed-forward neural network} \\
        % FPGA  & Field-programmable gate array \\
        GPU   & \link{graphics processing unit}{Graphics processing unit} \\
        IP    & \link{sec:1:IP_OOP}{In-plane} \\
        KMC   & \link{kinetic Monte Carlo}{Kinetic Monte Carlo} \\
        LIPO  & \link{LIPO}{Lipschitz Optimisation} \\
        LMC   & \link{memory capacity}{Linear memory capacity} \\
        LSM   & \link{liquid state machine}{Liquid state machine} \\
        MC    & \link{magnetostatic coupling}{Magnetostatic coupling}\footnote{Not to be confused with ``Monte Carlo'', which is always written in full, unless it is part of a longer abbreviation like KMC, MCMC or MCS.} \\
        MCMC  & \xref{Markov chain Monte Carlo} \\
        MCS   & \xref{Monte Carlo sweep} (i.e., 1 Monte Carlo step per magnet) \\
        MEMS  & \link{tensegrity}{Micro-electronic mechanical systems} \\
        MFM   & \link{magnetic force microscopy}{Magnetic force microscopy} \\
        MG    & \link{Mackey-Glass oscillator}{Mackey-Glass} \\
        MOKE  & \link{magneto-optical Kerr effect}{Magneto-optical Kerr effect} \\
        MRAM  & \link{spin-transfer torque}{Magnetic random-access memory} \\
        MS    & \link{magnetostatic interaction}{Magnetostatic} \\
        MSE   & \link{mean squared error}{Mean squared error} \\
        MTJ   & \link{spin-torque oscillator}{Magnetic tunnel junction} \\
        NL    & \link{non-linearity}{Non-linearity} \\
        NN    & \link{local antiferromagnetic parameter}{Nearest neighbour}\footnote{Second-nearest neighbours are abbreviated as 2NN; further neighbours as 3NN, etc.} \\
        % NTNU  & Norges teknisk-naturvitenskapelige universitet \\
        NV    & \link{nitrogen-vacancy microscopy}{Nitrogen-vacancy} \\
        OLS   & \link{ordinary least-squares}{Ordinary least-squares} \\
        OOP   & \link{sec:1:IP_OOP}{Out-of-plane} \\
        PBC   & \link{periodic boundary conditions}{Periodic boundary conditions} \\
        PC    & \link{parity check}{Parity check} \\
        PEEM  & \link{photo-emission electron microscopy}{Photo-emission electron microscopy} \\
        PMA   & \link{perpendicular magnetic anisotropy}{Perpendicular magnetic anisotropy} \\
        PRC   & \link{physical reservoir computing}{Physical reservoir computing} \\
        % PSI   & Paul Scherrer Institute \\
        RAM   & \link{von Neumann architecture}{Random-access memory} \\
        RC    & \link{reservoir computing}{Reservoir computing} \\
        RKKY  & \link{RKKY interaction}{Ruderman-Kittel-Kasuya-Yosida interaction} \\
        RNG   & \link{random number generator}{Random number generator} \\
        RNN   & \link{recurrent neural network}{Recurrent neural network} \\
        %RTRL  & Real-time recurrent learning \\
        SJG   & \link{stratified jittered grid}{Stratified jittered grid} \\
        SOT   & \link{spin--orbit torque}{Spin--orbit torque} \\
        STO   & \link{spin-torque oscillator}{Spin-torque oscillator} \\
        STT   & \link{spin-transfer torque}{Spin-transfer torque} \\
        STXM  & \link{scanning transmission X-ray microscopy}{Scanning transmission X-ray microscopy} \\
        %VLSI  & Very large-scale integration \\
        XMCD  & \xref{X-ray magnetic circular dichroism} \\
        %XOR   & Exclusive or \\
        YIG   & \link{spintronics}{Yttrium iron garnet} \\
    \end{longtable}
}
\cleardoublepage
\printindex\cleardoublepage
\chapter{Publications}
\renewcommand{\bibitem}[1]{\item}

\paragraph*{Papers}
\begin{itemize}
    \bibitem{MAES-24}
    \textbf{J.~Maes}, D.~De~Gusem, I.~Lateur, J.~Leliaert, A.~Kurenkov, and B.~Van~Waeyenberge.
    \newblock The design, verification, and applications of Hotspice: a Monte Carlo simulator for artificial spin ice.
    \newblock \emph{ArXiV}, arXiv:\penalty0 2409.05580, 2024.

    \bibitem{KUR-24}
    A.~Kurenkov, \textbf{J.~Maes}, A.~Pac, G.~M.~Macauley, B.~Van~Waeyenberge, A.~Hrabec, and L.~J.~Heyderman.
    \newblock Perpendicular-anisotropy artificial spin ice with spontaneous ordering: a platform for neuromorphic computing with flexible timescales.
    \newblock \emph{ArXiV}, arXiv:\penalty0 2408.12182, 2024.

    \bibitem{MOR-24}
    L.~Moreels, I.~Lateur, D.~De~Gusem, J.~Mulkers, \textbf{J.~Maes}, M.~V.~Milosevic, J.~Leliaert, and B.~Van~Waeyenberge.
    \newblock mumax+: extensible GPU-accelerated micromagnetics and beyond.
    \newblock \emph{ArXiV}, arXiv:\penalty0 2411.18194, 2024.
\end{itemize}

\paragraph*{Presentations}
\begin{itemize}
    \bibitem{SPINENGINE-24}
    \textbf{J.~Maes}, A.~Kurenkov, A.~Pac, G.~M.~Mcauley, B.~Van~Waeyenberge, A.~Hrabec and L.~J.~Heyderman.
    \newblock Reservoir Computing in Thermally Active Perpendicular-Anisotropy Artificial Spin Ice.
    \newblock \emph{Workshop on Magnetic Neuromorphic Computing 2024}, Trondheim (Norway), June 17-18, 2024.
\end{itemize}

\paragraph*{Other}
\begin{itemize}
    \bibitem{FASI-23}
    \textbf{J.~Maes} and B.~Van~Waeyenberge.
    \newblock Modelling of reservoir computing in out-of-plane artificial spin ice.
    \newblock \emph{Frontiers in Artificial Spin Ice 2023}, Lenzburg (Switzerland), June 14-16, 2023 (Poster).

    \bibitem{JEMS-23}
    \textbf{J.~Maes} and B.~Van~Waeyenberge.
    \newblock Modelling of reservoir computing in out-of-plane artificial spin ice.
    \newblock \emph{13th Joint European Magnetic Symposium}, Madrid (Spain), August 27-31, 2023 (Poster).

    \bibitem{MMM-23}
    \textbf{J.~Maes}, A.~Kurenkov and B.~Van~Waeyenberge.
    \newblock Modelling of Reservoir Computing in Out-of-Plane Artificial Spin Ice.
    \newblock \emph{68th Annual Conference on Magnetism and Magnetic Materials}, Dallas, Texas (USA), October 30 - November 3, 2023 (Poster).

    \bibitem{EAI-23}
    \textbf{J.~Maes}, A.~Kurenkov and B.~Van~Waeyenberge.
    \newblock Modelling of Reservoir Computing in Out-of-Plane Artificial Spin Ice.
    \newblock \emph{Emergent Algorithmic Intelligence Workshop 2023 `Bridging Disciplines: Novel AI-driven interdisciplinary techniques and methods'}, Mainz, Germany, December 6-8, 2023 (Poster).
\end{itemize}


\end{document}
